% Section 3, part 4 of 7: lifting, the scalar equation, proof of Theorem U, the leading term, the Z/5-degree, graded gauges.

\begin{lem}[Lifting]
\label{lem:gau-lift}
Let $\tilde A$ have the one-pole shape, let $X\in\mathfrak{gl}_5(\Bbbk\fps u)$ have shift $m+1$ with $0\leq m\leq4$, and suppose $\operatorname{pr}(X)=0$. Then there is $\mathcal Z\in\mathfrak{gl}_5(\Bbbk\fps u)$ of shift $m$ with $[\tilde A,\mathcal Z]=X$, and every solution $\mathcal Z'\in\mathfrak{gl}_5(\Bbbk\fps u)$ of $[\tilde A,\mathcal Z']=X$ is $\mathcal Z+w$ with $w\in\Bbbk\fps u[\mathcal Y]$. If $\tilde A$ and $X$ have rational coefficients, $\mathcal Z$ can be chosen rational.
\end{lem}

\begin{proof}
Write $X=\sum_a\xi_aE_{a,a+m+1}$ and look for $\mathcal Z=\sum_a\rho_a\Bprod^{(m)}_aE_{a,a+m}$ with $\rho_a\in\Bbbk\lser u$ (local letters). By \cref{lem:gau-comm}~(b) and $\tilde b_a\Bprod^{(m)}_{a+1}=\Bprod^{(m)}_a\tilde b_{a+m}=\Bprod^{(m+1)}_a$,
\[
[\tilde A,\mathcal Z]=\sum_a(\rho_{a+1}-\rho_a)\Bprod^{(m+1)}_aE_{a,a+m+1},
\]
so the equation is $\rho_{a+1}-\rho_a=\sigma_a:=\xi_a/\Bprod^{(m+1)}_a$ for $a\in\bbZ/5$. It is solvable if and only if $\sum_a\sigma_a=\lambda_{m+1}(X)=0$, i.e.\ if and only if $\operatorname{pr}(X)=0$ (\cref{lem:gau-comm}~(e)). Each $\sigma_a$ is integral, and $\sigma_a(0)=0$ whenever the window $\{a,\dots,a+m\}$ contains the pole index $1$, because then $\Bprod^{(m+1)}_a$ is $u^{-1}$ times a unit; for $m+1=5$ every window contains $1$. Put $\rho_1:=0$ and $\rho_{a+1}:=\rho_a+\sigma_a$ for $a=1,2,3,4$; the relation for $a=0$ then holds because $\sum_a\sigma_a=0$. All $\rho_a$ are integral. The entry $\rho_a\Bprod^{(m)}_a$ is integral if the window $\{a,\dots,a+m-1\}$ avoids $1$, and if it contains $1$ (that is, for the $m$ indices $a=2-m,\dots,1$) it is integral provided $\rho_a(0)=0$. For $a=1$ this holds by the choice $\rho_1=0$; for $a=2-m,\dots,0$ one has $\rho_a=\rho_1-\sum_{c=a}^0\sigma_c$, and every window $\{c,\dots,c+m\}$ with $2-m\leq c\leq0$ contains $1$, so $\rho_a(0)=0$. Hence $\mathcal Z$ is integral. Two integral solutions differ by an element of $\ker(\operatorname{ad}\tilde A)\cap\mathfrak{gl}_5(\Bbbk\fps u)=\Bbbk\fps u[\mathcal Y]$ (\cref{lem:gau-comm}~(c), (f)). The construction uses only field operations on the data.
\end{proof}

\begin{lem}[The scalar equation]
\label{lem:gau-scal}
Let $A$, $\tilde A$ have the one-pole shape and satisfy (EP).
\begin{enumerate}
\item[(a)] For $w=\sum_{j=0}^4w_j\mathcal Y^j$ with $w_j\in\Bbbk\lser u$,
\[
\operatorname{pr}\bigl(w'+\operatorname{dlog}\Theta_U\,w\bigr)=\sum_{j=0}^4\bigl(w_j'+\chi^U_jw_j\bigr)\mathcal Y^j,\qquad \chi^U_j:=-\frac j5\,\frac{\Bprod'}{\Bprod}+\frac15\tr(\operatorname{dlog}\Theta_U).
\]
In particular $\operatorname{pr}((\mathcal Y^j)')=-\frac j5(\Bprod'/\Bprod)\mathcal Y^j$ and $\operatorname{pr}(\operatorname{dlog}\Theta_U\,\mathcal Y^j)=\frac15\tr(\operatorname{dlog}\Theta_U)\,\mathcal Y^j$. Moreover $\chi^U_j=\frac j{5u}+\chi^\circ_j$ with $\chi^\circ_j:=-\frac j5(\Bprod^\circ)'/\Bprod^\circ+\frac15\tr(\operatorname{dlog}\Theta_U)\in\Bbbk\fps u$, where $\Bprod^\circ:=u\Bprod\in\Bbbk\fps u^\times$.
\item[(b)] For $Z\in\mathfrak{gl}_5(\Bbbk\fps u)$ the equation $\operatorname{pr}(w'+\operatorname{dlog}\Theta_U\,w)=-\operatorname{pr}(Z)$ has solutions $w\in\Bbbk\fps u[\mathcal Y]$; their components $w_1,\dots,w_4$ are unique, and $w_0$ is unique up to $w_0\mapsto w_0+\cn\,\mathfrak d$, $\cn\in\Bbbk$.
\item[(c)] For $Z=0$ the solutions in $\Bbbk\fps u[\mathcal Y]$ are exactly the $\cn\,\mathfrak d\,1$, $\cn\in\Bbbk$.
\item[(d)] If the data are rational, so are the solutions for $\cn\in\bbQ$; $\mathfrak d\in1+u\bbQ\fps u$.
\end{enumerate}
\end{lem}

\begin{proof}
(a) $(\mathcal Y^j)'=\sum_a((\Bprod^{(j)}_a)^{-1})'E_{a+j,a}$ has shift $5-j$; written as $\sum_{a'}\xi_{a'}E_{a',a'+5-j}$ it has $a'=a+j$ and $\xi_{a'}=((\Bprod^{(j)}_{a'-j})^{-1})'$. By \cref{lem:gau-comm}~(e) with $m=5-j$ (and $m=5$ for $j=0$),
\[
\operatorname{pr}((\mathcal Y^j)')=\frac15\Bigl[\sum_a\bigl((\Bprod^{(j)}_a)^{-1}\bigr)'\Bprod^{(j)}_a\Bigr]\mathcal Y^j=-\frac15\Bigl[\sum_a\frac{(\Bprod^{(j)}_a)'}{\Bprod^{(j)}_a}\Bigr]\mathcal Y^j .
\]
Here $(\Bprod^{(j)}_a)'/\Bprod^{(j)}_a=\sum_{i=0}^{j-1}\tilde b_{a+i}'/\tilde b_{a+i}$, and as $a$ runs over $\bbZ/5$ each index $c$ occurs for exactly the $j$ pairs $(a,i)=(c-i,i)$; so the sum is $j\sum_c\tilde b_c'/\tilde b_c=j\Bprod'/\Bprod$. Next, $\operatorname{dlog}\Theta_U\,\mathcal Y^j=\sum_a\frac{(d^U_{a+j})'}{d^U_{a+j}}(\Bprod^{(j)}_a)^{-1}E_{a+j,a}$, so now $\xi_{a'}=\frac{(d^U_{a'})'}{d^U_{a'}}(\Bprod^{(j)}_{a'-j})^{-1}$, and \cref{lem:gau-comm}~(e) gives $\frac15\bigl[\sum_{a'}(d^U_{a'})'/d^U_{a'}\bigr]\mathcal Y^j=\frac15\tr(\operatorname{dlog}\Theta_U)\mathcal Y^j$. As $\operatorname{pr}$ is $\Bbbk\lser u$-linear and fixes $\Bbbk\lser u[\mathcal Y]$, $\operatorname{pr}(w'+\operatorname{dlog}\Theta_Uw)=\sum_j\operatorname{pr}\bigl(w_j'\mathcal Y^j+w_j(\mathcal Y^j)'+w_j\operatorname{dlog}\Theta_U\mathcal Y^j\bigr)=\sum_j(w_j'+\chi^U_jw_j)\mathcal Y^j$. Finally $\Bprod=\Bprod^\circ/u$ with $\Bprod^\circ$ a unit by the one-pole shape, so $\Bprod'/\Bprod=(\Bprod^\circ)'/\Bprod^\circ-u^{-1}$, and $\tr(\operatorname{dlog}\Theta_U)=\sum_c(d^U_c)'/d^U_c$ is integral because the $d^U_c$ are units.

(b) By \cref{lem:gau-comm}~(f), $\operatorname{pr}(Z)=\sum_j\phi^Z_j\mathcal Y^j$ with $u\phi^Z_j\in\Bbbk\fps u$ for $j\geq1$ and $\phi^Z_0\in\Bbbk\fps u$ (local letters). By (a) and the linear independence of the $\mathcal Y^j$ the equation splits into $w_j'+\chi^U_jw_j=-\phi^Z_j$, $j=0,\dots,4$. For $j=1,\dots,4$ multiply by $u$: $uw_j'+(\frac j5+u\chi^\circ_j)w_j=-u\phi^Z_j$. For $w_j=\sum_nw_{j,n}u^n$ the coefficient of $u^n$ reads
\[
\Bigl(n+\frac j5\Bigr)w_{j,n}=-\coef{u^n}(u\phi^Z_j)-\sum_{l=0}^{n-1}\coef{u^{n-1-l}}\chi^\circ_j\cdot w_{j,l},
\]
and $n+j/5$ is a non-zero rational number for every $n\geq0$: so there is exactly one solution $w_j\in\Bbbk\fps u$. For $j=0$, $\chi^U_0=\frac15\tr(\operatorname{dlog}\Theta_U)=-\mathfrak d'/\mathfrak d$, so the equation is $\mathfrak d\,(w_0/\mathfrak d)'=-\phi^Z_0$; its solutions in $\Bbbk\fps u$ are $w_0=\mathfrak d\bigl(\cn-\int_0^u\phi^Z_0/\mathfrak d\,\dd u\bigr)$, $\cn\in\Bbbk$, where $\int_0^u$ is the formal antiderivative without constant term, and the difference $e$ of two solutions satisfies $(e/\mathfrak d)'=0$, so $e\in\Bbbk\mathfrak d$.

(c) is (b) with $\phi^Z=0$: the recursion gives $w_j=0$ for $j\geq1$, and $w_0=\cn\mathfrak d$. (d) The pivots $n+j/5$ are rational, $\det\Theta_U/\det\Theta_U(0)$ lies in $1+u\bbQ\fps u$ and its binomial series $\mathfrak d$ has rational coefficients.
\end{proof}

\begin{proof}[Proof of \cref{thm:gau-U}~(a)]
\emph{Reduction.} Write $P=\Theta_U\check P$; $\check P$ (a local letter) has coefficients in $\mathfrak{gl}_5(\Bbbk\fps u)$ if and only if $P$ has, since $\Theta_U$ is invertible over $\Bbbk\fps u$. As $A\Theta_U=\Theta_U\tilde A$, we have $AP-P\tilde A=\Theta_U(\tilde A\check P-\check P\tilde A)$ and $zP'=z\Theta_U'\check P+z\Theta_U\check P'$, so the gauge equation is equivalent to $z\check P'=[\tilde A,\check P]-z\operatorname{dlog}\Theta_U\,\check P$, i.e.\ to
\begin{equation}
\label{eq:gau-reduced}
[\tilde A,\check P_0]=0,\qquad [\tilde A,\check P_{k+1}]=\check P_k'+\operatorname{dlog}\Theta_U\,\check P_k\quad(k\geq0).
\end{equation}
\emph{Construction.} We construct $\check P_k=\mathcal Z_k+w^{(k)}$ with $\mathcal Z_k\in\mathfrak{gl}_5(\Bbbk\fps u)$ and $w^{(k)}\in\Bbbk\fps u[\mathcal Y]$ such that, for every $k\geq0$, (I$_k$) $[\tilde A,\check P_k]=\check P_{k-1}'+\operatorname{dlog}\Theta_U\check P_{k-1}$ (with $\check P_{-1}:=0$) and (S$_k$) $\operatorname{pr}(\check P_k'+\operatorname{dlog}\Theta_U\check P_k)=0$. For $k=0$ put $\mathcal Z_0:=0$ and $w^{(0)}:=\mathfrak d\,1$: (I$_0$) holds because $w^{(0)}$ commutes with $\tilde A$, and (S$_0$) is \cref{lem:gau-scal}~(c). Given $\check P_k$, the matrix $X_{k+1}:=\check P_k'+\operatorname{dlog}\Theta_U\check P_k$ is integral with $\operatorname{pr}(X_{k+1})=0$ by (S$_k$); since $\operatorname{pr}$ preserves the shift, \cref{lem:gau-lift} applied to each shift component gives an integral $\mathcal Z_{k+1}$ with $[\tilde A,\mathcal Z_{k+1}]=X_{k+1}$. Choose $w^{(k+1)}\in\Bbbk\fps u[\mathcal Y]$ by \cref{lem:gau-scal}~(b) with $Z:=\mathcal Z_{k+1}'+\operatorname{dlog}\Theta_U\mathcal Z_{k+1}$ and any value of the free constant. Then (I$_{k+1}$) holds since $[\tilde A,w^{(k+1)}]=0$, and (S$_{k+1}$) holds since $\operatorname{pr}(\check P_{k+1}'+\operatorname{dlog}\Theta_U\check P_{k+1})=\operatorname{pr}(Z)+\operatorname{pr}\bigl((w^{(k+1)})'+\operatorname{dlog}\Theta_Uw^{(k+1)}\bigr)=0$. By \eqref{eq:gau-reduced}, $P^\natural:=\Theta_U\sum_kz^k\check P_k$ lies in $\Gau(A,\tilde A)$, and $P^\natural_0=\mathfrak d\,\Theta_U$. The only choices are the free constant of \cref{lem:gau-scal}~(b) at each order and the normalisation $\rho_1=0$ inside \cref{lem:gau-lift}, which the next choice of $w^{(k+1)}$ absorbs.
\end{proof}

\begin{lem}[The leading term of a solution of the gauge equation]
\label{lem:gau-lead}
Let $A$, $\tilde A$ have the one-pole shape and satisfy (EP), let $P\in\Gau(A,\tilde A)$ and $n\geq0$ with $P_0=\dots=P_{n-1}=0$. Then $P_n=\cn\,\mathfrak d\,\Theta_U$ for a unique $\cn\in\Bbbk$, namely $\cn=(P_n)_{00}(0)$.
\end{lem}

\begin{proof}
Put $\check P:=\Theta_U^{-1}P$, so $\check P_l=0$ for $l<n$. The coefficient of $z^n$ in the reduced equation \eqref{eq:gau-reduced} is $[\tilde A,\check P_n]=0$ (for $n=0$ this is the first equation), so $\check P_n\in\ker(\operatorname{ad}\tilde A)\cap\mathfrak{gl}_5(\Bbbk\fps u)=\Bbbk\fps u[\mathcal Y]$ (\cref{lem:gau-comm}~(c), (f)). The coefficient of $z^{n+1}$ says that $\check P_n'+\operatorname{dlog}\Theta_U\check P_n=[\tilde A,\check P_{n+1}]$ lies in $\operatorname{im}(\operatorname{ad}\tilde A)$, so $\operatorname{pr}(\check P_n'+\operatorname{dlog}\Theta_U\check P_n)=0$. By \cref{lem:gau-scal}~(c), $\check P_n=\cn\,\mathfrak d\,1$. As $d^U_0=1$ and $\mathfrak d(0)=1$, $\cn=(P_n)_{00}(0)$.
\end{proof}

\begin{proof}[Proof of \cref{thm:gau-U}~(b) and (d)]
(b) Let $P^\natural\in\Gau(A,\tilde A)$ with $P^\natural_0=\mathfrak d\,\Theta_U$; the gauge of (a) is one. \emph{Injectivity:} if $\cn\neq0$ has order $n$ in $z$, then $\coef{z^n}(P^\natural\cn)=\cn_nP^\natural_0=\cn_n\mathfrak d\Theta_U\neq0$. \emph{Surjectivity:} let $P\in\Gau(A,\tilde A)$. We define $\cn_0,\cn_1,\dots\in\Bbbk$ recursively. Suppose $R^{(n)}:=P-P^\natural(\cn_0+\cn_1z+\dots+\cn_{n-1}z^{n-1})$ is $0$ modulo $z^n$ (for $n=0$ this is empty); $R^{(n)}\in\Gau(A,\tilde A)$ because $\Gau(A,\tilde A)$ is a $\Bbbk\fps z$-module. By \cref{lem:gau-lead}, $\coef{z^n}R^{(n)}=\cn_n\mathfrak d\Theta_U$ for a unique $\cn_n\in\Bbbk$. Then $R^{(n+1)}=R^{(n)}-\cn_nz^nP^\natural$ lies in $\Gau(A,\tilde A)$ and $\coef{z^n}R^{(n+1)}=\cn_n\mathfrak d\Theta_U-\cn_nP^\natural_0=0$, so $R^{(n+1)}\equiv0$ modulo $z^{n+1}$. Put $\cn:=\sum_n\cn_nz^n$. For every $l$ and every $n>l$, $\coef{z^l}(P-P^\natural\cn)=\coef{z^l}R^{(n)}-\coef{z^l}P^\natural(\cn-\sum_{m<n}\cn_mz^m)=0$. Hence $P=P^\natural\cn$. Its $z^0$-term is $\cn(0)\mathfrak d\Theta_U$, invertible if and only if $\cn(0)\neq0$. If $P=P^\natural\cn_P$ is a regular gauge, every regular gauge $P^\natural\cn'$ equals $P\cdot(\cn_P^{-1}\cn')$ with $\cn_P^{-1}\cn'\in\Bbbk\fps z^\times$, and conversely. For the pair $(\tilde A,\tilde A)$ one has $\Theta_U=1$ and $\mathfrak d=1$, and $1\in\Gau(\tilde A,\tilde A)$; so $\Gau(\tilde A,\tilde A)=\Bbbk\fps z\cdot1$.

(d) Every object used in the proof of (a) is obtained from the $b_a$, $\tilde b_a$ by field operations in $\bbQ\lser u$, the derivative and the integration $\int_0^u$ of power series: $\Theta_U$, $\operatorname{dlog}\Theta_U$, $\mathcal Y$, $\Bprod^{(m)}_a$, $\Bprod^\circ$, $\operatorname{pr}$ (\cref{lem:gau-comm}~(e)), $\mathfrak d$, the lift of \cref{lem:gau-lift} and the solutions of \cref{lem:gau-scal}~(b) (\cref{lem:gau-scal}~(d)). With every free constant chosen in $\bbQ$, the construction therefore runs inside $\mathfrak{gl}_5(\bbQ\fps u)\fps z$. For a field $\Bbbk\supseteq\bbQ$ the proofs of \cref{lem:gau-comm,lem:gau-lift,lem:gau-scal,lem:gau-lead} hold verbatim over $\Bbbk$, so the solutions over $\Bbbk$ are $P^\natural\Bbbk\fps z$ by (b). If $P=P^\natural\cn$ has rational coefficients, the recursion of the surjectivity proof computes $\cn_n=(\coef{z^n}R^{(n)})_{00}(0)$ from rational data, so $\cn\in\bbQ\fps z$.
\end{proof}

\begin{lem}[The $\bbZ/5$-degree of the gauge equation]
\label{lem:gau-degree}
For $P=\sum_kz^kP_k$ with $P_k\in\mathfrak{gl}_5(\Bbbk\lser u)$ and $\delta\in\bbZ/5$ put $P^{[\delta]}:=\sum_kz^k\sum_{(a,b):\,k+b-a\equiv\delta}(P_k)_{ab}E_{ab}$, so that $P=\sum_\delta P^{[\delta]}$, and $P$ is graded if and only if $P=P^{[0]}$. Let $\operatorname{ge}(P):=zP'-AP+P\tilde A$. Then $\operatorname{ge}(P^{[\delta]})=\operatorname{ge}(P)^{[\delta+1]}$; in particular every degree part of a solution of the gauge equation is a solution.
\end{lem}

\begin{proof}
Give $z^kf(u)E_{ab}$ the degree $k+b-a$ modulo $5$. The derivative in $u$ does not change it and the factor $z$ adds $1$, so $zP'$ raises the degree of each homogeneous term by $1$. $A\,(fz^kE_{cd})=b_{c-1}fz^kE_{c-1,d}$ has degree $k+d-c+1$, and $(fz^kE_{cd})\tilde A=\tilde b_dfz^kE_{c,d+1}$ has degree $k+d+1-c$. So $\operatorname{ge}$ maps homogeneous series of degree $\delta$ to homogeneous series of degree $\delta+1$, and by additivity $\operatorname{ge}(P)^{[\delta+1]}=\operatorname{ge}(P^{[\delta]})$.
\end{proof}

\begin{prop}[Graded gauges]
\label{prop:gau-graded}
Let $A$, $\tilde A$ have the one-pole shape and satisfy (EP).
\begin{enumerate}
\item Graded regular gauges exist; they can be chosen rational if all $b_a$, $\tilde b_a$ are rational.
\item If $P$ is a graded regular gauge, the graded regular gauges are exactly the $P\cn(z)$ with $\cn\in\Bbbk\fps{z^5}^\times$.
\item For a graded regular gauge $P$, $P_{11}(0,z)=\sum_{m\geq0}(P_{5m})_{11}(0)\,z^{5m}$ is a unit of $\Bbbk\fps{z^5}$.
\item If $\zeta\in\Bbbk$, then $\operatorname{Rot}_iP=\sum_\delta\zeta^{-i\delta}P^{[\delta]}$; so $P$ is graded if and only if $\operatorname{Rot}_1P=P$, if and only if $\operatorname{Rot}_iP=P$ for all $i$, and $\frac15\sum_i\operatorname{Rot}_iP=P^{[0]}$. Moreover $A\,\zeta^{i(\mathrm{gr}+1)}=\zeta^i\zeta^{i(\mathrm{gr}+1)}A$, and the same for $\tilde A$; $\operatorname{Rot}_i$ maps solutions of $z\de_uY=AY$ to solutions, and $\operatorname{Rot}_i(PY)=(\operatorname{Rot}_iP)(\operatorname{Rot}_iY)$, $\operatorname{Rot}_i(Y\mathsf d(z))=\operatorname{Rot}_i(Y)\,\mathsf d(\zeta^{-i}z)$ for scalar series $\mathsf d$.
\end{enumerate}
\end{prop}

\begin{proof}
(i) Let $P^\natural$ be as in \cref{thm:gau-U}~(a), rational in the rational case by (d). Its $z^0$-term $\mathfrak d\Theta_U$ is diagonal, of degree $0$. By \cref{lem:gau-degree}, $P^{\mathrm{gr}}:=(P^\natural)^{[0]}$ solves the gauge equation, and $P^{\mathrm{gr}}_0=P^\natural_0$ is invertible. (ii) By \cref{thm:gau-U}~(b) the regular gauges are the $P\cn(z)$, $\cn\in\Bbbk\fps z^\times$. Write $\cn=\sum_{r\in\bbZ/5}\cn^{(r)}$ with $\cn^{(r)}:=\sum_{n\equiv r}\cn_nz^n$. As $P$ has degree $0$ and $z^n$ adds $n$ to the degree, $(P\cn)^{[r]}=P\cn^{(r)}$. So $P\cn$ is graded if and only if $P\cn^{(r)}=0$ for $r\neq0$, if and only if $\cn^{(r)}=0$ for $r\neq0$, because $\cn\mapsto P\cn$ is injective ($P=P^\natural\cn_P$ with $\cn_P\neq0$, and $\Bbbk\fps z$ is a domain); that is, if and only if $\cn\in\Bbbk\fps{z^5}$. (iii) For graded $P$ the entry $(P_k)_{11}$ vanishes unless $1\equiv1+k$, so $P_{11}(0,z)\in\Bbbk\fps{z^5}$, and its constant term is $(P_0)_{11}(0)=\cn(0)\mathfrak d(0)d^U_1(0)\neq0$ by \cref{thm:gau-U}~(b). (iv) The entry $(a,b)$ of $\coef{z^k}\operatorname{Rot}_iP$ is $\zeta^{i(a+1)-ik-i(b+1)}(P_k)_{ab}=\zeta^{-i(k+b-a)}(P_k)_{ab}$, which gives the first assertion; the equivalences follow because the decomposition $P=\sum_\delta P^{[\delta]}$ is direct, $\zeta^{-i\delta}\neq1$ for $i,\delta\not\equiv0$, and $\sum_i\zeta^{-i\delta}=5\,\delta_{\delta,0}$. The entry $(a,a+1)$ of $\zeta^{i(\mathrm{gr}+1)}A\zeta^{-i(\mathrm{gr}+1)}$ is $\zeta^{i(a+1)}b_a\zeta^{-i(a+2)}=\zeta^{-i}b_a$, which is the commutation rule. Hence, for a solution $Y$, $z\de_u(\operatorname{Rot}_iY)=\zeta^i\zeta^{i(\mathrm{gr}+1)}\bigl[(\zeta^{-i}z)\de_uY\bigr](u,\zeta^{-i}z)=\zeta^i\zeta^{i(\mathrm{gr}+1)}A\,Y(u,\zeta^{-i}z)=A\operatorname{Rot}_iY$. The two product rules are immediate from the definitions.
\end{proof}

\begin{rmk}[What the proof of Theorem~U uses; the role of the pole]
\label{rmk:gau-Uremarks}
The proof of \cref{thm:gau-U} uses only the one-pole shape, (EP) and the characteristic $0$ (in $\lambda_m(\tilde A^m)=5\neq0$, in the pivots $n+j/5\neq0$, and for $\mathfrak d$ and the constants of $\dd/\dd u$). The pole of $\tilde b_1$ enters exactly through the integrality statements, \cref{lem:gau-comm}~(f) and the window argument of \cref{lem:gau-lift}, and through the term $j/(5u)$ of $\chi^U_j$. For two systems without a pole (all ten weights units, with (EP)) the integrality statements hold trivially and $\chi^U_j\in\Bbbk\fps u$ for every $j$, so in the proof of \cref{lem:gau-scal}~(b) every equation $w_j'+\chi^U_jw_j=-\phi^Z_j$, $j=0,\dots,4$, has the solutions $\mathcal E_j\bigl(\mathsf a_j-\int_0^u\phi^Z_j/\mathcal E_j\,\dd u\bigr)$, $\mathsf a_j\in\Bbbk$, with $\mathcal E_j:=\exp(-\int_0^u\chi^U_j\,\dd u)$ (local letters). Then the proof of \cref{lem:gau-lead} gives the leading terms $\Theta_U\sum_j\mathsf a_j\mathcal E_j\mathcal Y^j$, the construction of \cref{thm:gau-U}~(a) started from $w^{(0)}:=\mathcal E_j\mathcal Y^j$ gives a solution $P^{(j)}$ for each $j$, and the proof of \cref{thm:gau-U}~(b), with the five leading terms in place of one, shows that $P^{(0)},\dots,P^{(4)}$ form a basis of $\Gau(A,\tilde A)$ over $\Bbbk\fps z$. So the solution module has rank $5$, and $P_{11}(0,z)=1$ no longer fixes a gauge (see \cref{sec:evi-gauge} for a computed example). This is not used.
\end{rmk}

% Section 3, part 5 of 7: (R1), the diagonal, the vanishing solutions, the gauge identity (G), the canonical gauge, symplecticity.

%-------------------------------------------------------------------------------
\subsection{The residue relation, the diagonal, the vanishing solutions and (G)}
\label{sec:gau-van}

From now on, unless said otherwise, $A$ is the quintic system of \cref{prop:gau-quintic} in a conifold frame and $\tilde A$ is the model pulled back along $v=v(u)$ (\cref{prop:gau-equalprod}). Both are self-dual cyclic systems with one simple pole over $\bbQ\fps u$, with (EP), $r=\tilde r=-1$ and $b_a(0)=\tilde b_a(0)$ for $a\neq1$, so \cref{thm:gau-U} applies over $\bbQ$.

\begin{lem}[The residue relation (R1)]
\label{lem:gau-R1}
Let $A,\tilde A$ be cyclic systems with one simple pole over $\Bbbk\fps u$, with residues $r,\tilde r$, and let $P\in\Gau(A,\tilde A)$; put $P(0,z):=\sum_kz^kP_k(0)\in\mathfrak{gl}_5(\Bbbk\fps z)$. Then $rE_{12}P(0,z)=\tilde r\,P(0,z)E_{12}$: column $1$ of $P(0,z)$ is $P_{11}(0,z)\mathbf e_1$, row $2$ is $P_{22}(0,z)\mathbf e_2^T$, and $rP_{22}(0,z)=\tilde rP_{11}(0,z)$. If $r=\tilde r$, then $P_{22}(0,z)=P_{11}(0,z)$.
\end{lem}

\begin{proof}
$A=u^{-1}rE_{12}+(\text{power series in }u)$, and likewise for $\tilde A$. The term $zP'$ has no $u^{-1}$, so the coefficient of $u^{-1}$ in $AP-P\tilde A$ vanishes: $rE_{12}P(0,z)=\tilde rP(0,z)E_{12}$. The entry $(a,b)$ reads $r\,\delta_{a1}P_{2b}(0,z)=\tilde r\,P_{a1}(0,z)\delta_{b2}$; for $a=1$, $b\neq2$ it gives $P_{2b}(0,z)=0$, for $a\neq1$, $b=2$ it gives $P_{a1}(0,z)=0$, and for $(a,b)=(1,2)$ it gives $rP_{22}(0,z)=\tilde rP_{11}(0,z)$.
\end{proof}

\begin{lem}[The diagonal of the gauge]
\label{lem:gau-diag}
Let $A,\tilde A$ be cyclic systems with one simple pole over $\Bbbk\fps u$ with (EP), and $P$ a regular gauge. Then $P_0=\diag(p_{0,0},\dots,p_{0,4})$, and $\prod_ap_{0,a}=\cn(0)^5\det\Theta_U(0)=p_{0,1}(0)^5\prod_a\rho^U_a$ is a constant, where $\cn(0)\in\Bbbk^\times$ is as in \cref{thm:gau-U}~(b) and $\rho^U_a:=d^U_a(0)/d^U_1(0)$, i.e.\ $\rho^U_1=1$, $\rho^U_a=\prod_{c=1}^{a-1}(\tilde b_c/b_c)(0)$ for $a=2,3,4$ and $\rho^U_0=\prod_{c=1}^4(\tilde b_c/b_c)(0)$, with $(\tilde b_1/b_1)(0):=\tilde r/r$. Assume the normalisation
\[
\text{(N0)}\qquad b_a(0)=\tilde b_a(0)\ \ (a\neq1),\qquad r=\tilde r .
\]
Then $p_{0,a}(0)=p_{0,1}(0)$ for all $a$; if moreover $p_{0,1}(0)=1$, as for the canonical gauge of \cref{def:gau-can}, then $p_{0,a}(0)=1$ for all $a$, $\prod_ap_{0,a}=1$, and $p_{0,0}^5=\prod_{c=0}^3(b_c/\tilde b_c)^{4-c}$. For the quintic system in a conifold frame and the pulled-back model, where (N0) holds, and a regular gauge with $p_{0,1}(0)=1$ (for instance $\Pc$), this gives
\[
p_{0,0}(u)=\tau(u)^{-1/4}/\varpi_0,\qquad\text{i.e.}\qquad \tau^{1/4}p_{0,0}=1/\varpi_0 .
\]
\end{lem}

\begin{proof}
By \cref{thm:gau-U}~(b), $P_0=\cn(0)\,\mathfrak d\,\Theta_U$ with $\cn(0)\neq0$; so $p_{0,a}=\cn(0)\mathfrak d\,d^U_a$ and $\prod_ap_{0,a}=\cn(0)^5\mathfrak d^5\det\Theta_U=\cn(0)^5\det\Theta_U(0)$, by the definition of $\mathfrak d$. At $u=0$, $p_{0,a}(0)=\cn(0)d^U_a(0)=p_{0,1}(0)\rho^U_a$. The recursion $d^U_{a+1}=d^U_a\tilde b_a/b_a$ at $u=0$, with $(\tilde b_1/b_1)(0)=\tilde r/r$ and, going round the cycle by (EP), $d^U_0=d^U_4\tilde b_4/b_4$, gives the formula for $\rho^U_a$. Under (N0) every ratio at $u=0$ is $1$, so $d^U_a(0)=1$ and $\rho^U_a=1$ for all $a$, and $p_{0,a}(0)=\cn(0)$. If $p_{0,1}(0)=1$, then $\cn(0)=1$, all $p_{0,a}(0)=1$, $\prod_ap_{0,a}=\det\Theta_U(0)=1$, and $p_{0,0}=\mathfrak d=(\det\Theta_U)^{-1/5}$. Since $d^U_a=\prod_{c<a}\tilde b_c/b_c$, each $c\in\{0,\dots,3\}$ occurs in $4-c$ of the $d^U_a$, so $\det\Theta_U=\prod_{c=0}^3(\tilde b_c/b_c)^{4-c}$.

For the two systems of \cref{prop:gau-quintic,prop:gau-equalprod} the ratios $b_c/\tilde b_c$, $c=0,\dots,3$, are
\[
\frac{\tc'}{v'},\qquad \frac{v}{u\,v'\varpi_0^2(1-u)^3\tc'^2},\qquad\frac{\tc'}{v'},\qquad\frac{\varpi_0}{(1-u)v'},
\]
so $p_{0,0}^5=\tc'^{\,4-6+2}\,v'^{-10}\,v^3u^{-3}\,\varpi_0^{-6+1}(1-u)^{-9-1}=v'^{-10}v^3u^{-3}\varpi_0^{-5}(1-u)^{-10}$. By (EP), $v'^5=v/(u(1-u)^5)$ (\cref{prop:gau-equalprod}~(ii)), hence $v'^{-10}=u^2(1-u)^{10}v^{-2}$ and $p_{0,0}^5=v/(u\varpi_0^5)=\tau^{-5/4}\varpi_0^{-5}$. Both $p_{0,0}$ and $\tau^{-1/4}/\varpi_0$ lie in $1+u\bbQ\fps u$, where fifth roots are unique, so they are equal.
\end{proof}

\begin{defn}[Vanishing solutions]
\label{def:gau-van}
\begin{enumerate}
\item For $m\in\bbZ$ let $\IC_m(v,z):=\Res_{p=0}p^m\re^{-\Wv(p)/z}\dd p$ be the formal residue: in $\re^{-\Wv(p)/z}=\re^{-p^4/(4z)}\re^{-v/(pz)}$ the coefficient of each power of $v$ is a Laurent series in $p$ with coefficients in $\bbQ[z^{-1}]$, and $\Res_{p=0}$ takes its coefficient of $p^{-1}$. Explicitly
\[
\IC_m(v,z)=\sum_{i\geq0,\ m+4i+1\geq0}\frac{(-1/(4z))^i}{i!}\,\frac{(-v/z)^{m+4i+1}}{(m+4i+1)!}\in\bbQ[z^{-1}]\fps v .
\]
The \emph{vanishing solution of the model} is $\YC:=(\IC_0,-\IC_{-1},-\IC_3,\IC_2,\IC_1)$, i.e.\ $\YC_b=\varepsilon_b\IC_{\nu_b}$ (\cref{def:set-model}).
\item For a cyclic system $A=\sum_ab_aE_{a,a+1}$ with one simple pole over $\bbQ\fps u$ put $y^{(0)}:=-\mathbf e_1$, $y^{(n+1)}_a:=\int_0^ub_a\,y^{(n)}_{a+1}\,\dd u$ (the integrands are power series by \cref{lem:gau-vanq}) and $Y^{\mathrm C}_A:=\sum_{n\geq0}z^{-n}y^{(n)}$. For the quintic system, $\YCQ:=Y^{\mathrm C}_A$ is the \emph{vanishing solution of the quintic}.
\end{enumerate}
\end{defn}

\begin{lem}[The model's vanishing solution]
\label{lem:gau-vanmod}
(i) $z\,\de_v\IC_m=-\IC_{m-1}$; (ii) $mz\IC_{m-1}-\IC_{m+3}+v\IC_{m-2}=0$. Hence $z\,\dd\YC_b/\dd v=b^{\mathrm m}_b\YC_{b+1}$, $b^{\mathrm m}=(1,-v^{-1},1,-1,-1)$, and $\YC(0,z)=-\mathbf e_1$.
\end{lem}

\begin{proof}
(i) Differentiate the residue coefficientwise in $v$, using $\de_v\re^{-\Wv/z}=-(pz)^{-1}\re^{-\Wv/z}$. (ii) We have $\frac{\dd}{\dd p}\bigl(p^m\re^{-\Wv/z}\bigr)=\bigl(mp^{m-1}-z^{-1}(p^{m+3}-vp^{m-2})\bigr)\re^{-\Wv/z}$, since $\Wv'(p)=p^3-vp^{-2}$, and the residue of the derivative of a Laurent series vanishes. With $\YC=(\IC_0,-\IC_{-1},-\IC_3,\IC_2,\IC_1)$, (i) gives $z\de_v\YC_0=-\IC_{-1}=\YC_1$, $z\de_v\YC_2=\IC_2=\YC_3$, $z\de_v\YC_3=-\IC_1=-\YC_4$, $z\de_v\YC_4=-\IC_0=-\YC_0$, and $z\de_v\YC_1=\IC_{-2}=v^{-1}\IC_3=-v^{-1}\YC_2$ by (ii) at $m=0$. At $v=0$ only the terms with $m+4i+1=0$ survive; among $m\in\{0,-1,3,2,1\}$ this happens only for $m=-1$, $i=0$, with value $1$; so $\YC(0,z)=-\mathbf e_1$.
\end{proof}

\begin{lem}[Shape of $\YC$]
\label{lem:gau-vanshape}
For $n\geq0$ and $b\in\bbZ/5$, with $m:=\lfloor n/5\rfloor$,
\[
\coef{z^{-n}}\YC_b=\varepsilon_b\,(-1)^n\,\frac{v^{n-m}}{4^m\,m!\,(n-m)!}\quad\text{if }n\equiv\nu_b+1\pmod5,\qquad \coef{z^{-n}}\YC_b=0\ \text{otherwise}.
\]
Hence $\YC_b(v,z)=\varsigma^{\nu_b+1}y_b(z/\varsigma^4)$ with $y_b(t):=\YC_b(1,t)$ and $\varsigma=v^{1/5}$; the lowest terms are $\YC=(-v/z,\,-1,\,-v^4/(24z^4),\,-v^3/(6z^3),\,v^2/(2z^2))+\dots$; and at $v=v(u)$ the coefficient of $z^{-n}$ is $O(u^{n-m})$, so $\YC(v(u),z)\in\bbQ[z^{-1}]\fps u^5$ and $\YC(v(u),z)$ has no positive power of $z$.
\end{lem}

\begin{proof}
In \cref{def:gau-van}~(i) with $m'=\nu_b$, the term of index $i$ is a multiple of $v^{m'+4i+1}z^{-n}$ with $n=m'+5i+1$; as $\nu_b+1\in\{0,\dots,4\}$, this gives $n\equiv\nu_b+1$ and $i=\lfloor n/5\rfloor=m$, the exponent of $v$ is $n-m$, the sign is $(-1)^{i+m'+4i+1}=(-1)^n$, and the denominator is $4^mm!\,(n-m)!$. Since $v^{n-m}z^{-n}=\varsigma^{n-5m}(z/\varsigma^4)^{-n}$ and $n-5m=\nu_b+1$, the homogeneity follows. For $v=u+O(u^2)$ the coefficient of $z^{-n}$ is $O(u^{n-m})$ with $n-m\geq4n/5$, so each power of $u$ receives contributions from finitely many $n$.
\end{proof}

\begin{lem}[Uniqueness of regular solutions]
\label{lem:gau-uniq}
Let $B=u^{-1}B_{-1}+\sum_{j\geq0}u^jB_j$ with $B_j\in\mathfrak{gl}_5(\Bbbk)$ independent of $z$. If $Y\in\Bbbk\lser z\fps u^5$ satisfies $z\de_uY=BY$ and $Y(0)=0$, then $Y=0$.
\end{lem}

\begin{proof}
Write $Y=\sum_nu^nY^{(n)}$. The coefficient of $u^n$ in $z\de_uY=BY$ gives $\bigl((n+1)z-B_{-1}\bigr)Y^{(n+1)}=\sum_{j=0}^nB_jY^{(n-j)}$. The determinant $\det((n+1)z\,1-B_{-1})$ is a polynomial of degree $5$ in $z$ with leading coefficient $(n+1)^5\neq0$, hence a unit of $\Bbbk\lser z$ (for $B_{-1}=rE_{12}$ it is $((n+1)z)^5$). By induction all $Y^{(n)}$ vanish.
\end{proof}

A cyclic system with one simple pole is of this form ($B_{-1}=rE_{12}$), so its solutions in $\bbQ\lser z\fps u^5$ are determined by their values at $u=0$. By \cref{lem:gau-vanmod} and the chain rule, $\YC(v(u),z)$ is the solution of $z\de_uY=\tilde AY$ with value $-\mathbf e_1$, and by the next lemma $\YCQ$ is that of the quintic system; the vanishing solutions are characterised by this value, not by the absence of positive powers of $z$.

\begin{lem}[The quintic's vanishing solution]
\label{lem:gau-vanq}
Let $A=\sum_ab_aE_{a,a+1}$ be a cyclic system with one simple pole over $\bbQ\fps u$, with residue $r$.
\begin{enumerate}
\item[(a)] $y^{(n)}_a=0$ unless $n\equiv1-a\pmod5$, and then $\ord_uy^{(n)}_a=(4n+\nu_a+1)/5$ exactly. If $b_a(0)=\tilde b_a(0)$ for $a\neq1$ and $r=\tilde r$, where $\tilde A$ is the pulled-back model, the leading coefficient of $y^{(n)}_a$ equals that of $\coef{z^{-n}}\YC_a(v(u),z)$.
\item[(b)] Every integrand $b_ay^{(n)}_{a+1}$ is a power series: the pole of $b_1$ meets $y^{(n)}_2$, which is non-zero only for $n\equiv4\pmod5$ and then has order $(4n+4)/5\geq4$. In particular no logarithm arises.
\item[(c)] $Y^{\mathrm C}_A\in\bbQ[z^{-1}]\fps u^5$, the coefficient of $z^{-n}$ having order at least $4n/5$; it solves $z\de_uY=AY$, and $Y^{\mathrm C}_A(0)=-\mathbf e_1$.
\item[(d)] $\coef{z^{-1}}Y^{\mathrm C}_A=-\bigl(\int_0^ub_0\,\dd u\bigr)\mathbf e_0$ and $\coef{u^0}(Y^{\mathrm C}_A)_0=0$. For the quintic system in a conifold frame, $\coef{z^{-1}}\YCQ=-\tc\,\mathbf e_0$.
\end{enumerate}
\end{lem}

\begin{proof}
Put $e(n,a):=(4n+\nu_a+1)/5$; for $n\equiv1-a$ it is a non-negative integer, as $\nu_a\equiv-a$ and $\nu_a+1\geq0$. Write $b_a=b^\circ_au^{-\delta_{a1}}(1+O(u))$ with $b^\circ_a:=b_a(0)$ for $a\neq1$ and $b^\circ_1:=r$ (local letters). We show by induction on $n$ that for $n\equiv1-a$, $y^{(n)}_a=c^{(n)}_au^{e(n,a)}(1+O(u))$ with $c^{(n)}_a\neq0$, and $y^{(n)}_a=0$ otherwise. For $n=0$, $y^{(0)}=-\mathbf e_1$ and $e(0,1)=0$. The component $y^{(n+1)}_a$ is built from $y^{(n)}_{a+1}$, which vanishes unless $n\equiv1-(a+1)$, i.e.\ $n+1\equiv1-a$. In that case, using $\nu_{a+1}=\nu_a-1+5\delta_{a1}$, the integrand is $b_ay^{(n)}_{a+1}=b^\circ_ac^{(n)}_{a+1}u^{e(n+1,a)-1}(1+O(u))$ with $e(n+1,a)-1=(4n+\nu_a)/5\geq0$, a power series, and integration gives $c^{(n+1)}_a=b^\circ_ac^{(n)}_{a+1}/e(n+1,a)\neq0$. This proves the first part of (a) and (b); for $a=1$ the integrand involves $y^{(n)}_2$ with $n\equiv4$, of order $e(n,2)=(4n+4)/5\geq4$. The pulled-back model satisfies the same recursion for the leading coefficients with $\tilde b^\circ_a$ in place of $b^\circ_a$, so under the stated hypothesis the leading coefficients of $Y^{\mathrm C}_A$ and $Y^{\mathrm C}_{\tilde A}$ agree; and $Y^{\mathrm C}_{\tilde A}=\YC(v(u),z)$ by \cref{lem:gau-uniq}, both being solutions of $z\de_uY=\tilde AY$ with value $-\mathbf e_1$ (the first by (c) below, applied to $\tilde A$). (c) The order bound is $e(n,a)\geq4n/5$, so each power of $u$ receives finitely many $n$. The definition says $(y^{(n+1)})'=Ay^{(n)}$, hence $z\de_uY^{\mathrm C}_A=\sum_{n\geq1}z^{1-n}Ay^{(n-1)}=AY^{\mathrm C}_A$; and $y^{(n)}(0)=0$ for $n\geq1$. (d) $y^{(1)}_a=-\delta_{a0}\int_0^ub_0$, and $(Y^{\mathrm C}_A)_0(0)=y^{(0)}_0=0$. For the quintic, $b_0=\tc'$ and $\tc(0)=0$.
\end{proof}

In a conifold frame the hypothesis of (a) holds (\cref{prop:gau-equalprod}).

\begin{thm}[The gauge identity (G)]
\label{thm:gau-G}
For every $P\in\Gau(A,\tilde A)$, in particular for every regular gauge from the pulled-back model to the quintic system,
\begin{equation}
\label{eq:gau-G}
P(u,z)\,\YC(v(u),z)=P_{11}(0,z)\,\YCQ(u,z)\qquad\text{in }\bbQ\lser z\fps u^5 .
\end{equation}
\begin{enumerate}
\item[(a)] For a graded regular gauge $P$, $P_{11}(0,z)$ is a unit of $\bbQ\fps{z^5}$, $\Pc:=P\,P_{11}(0,z)^{-1}$ is a graded regular gauge, and $\Pc$ is the unique regular gauge $P'$ with $P'\,\YC(v(u),z)=\YCQ$.
\item[(b)] For a regular gauge $P$, $P\,\YC(v(u),z)$ has no positive power of $z$ if and only if $P_{11}(0,z)$ is a constant; the regular gauges with this property are the $\mu\Pc$, $\mu\in\bbQ^\times$.
\item[(c)] For any two cyclic systems $A,\tilde A$ with one simple pole over $\bbQ\fps u$ and every $P\in\Gau(A,\tilde A)$: $P\,Y^{\mathrm C}_{\tilde A}=P_{11}(0,z)\,Y^{\mathrm C}_A$.
\end{enumerate}
\end{thm}

\begin{proof}
We prove (c), of which \eqref{eq:gau-G} is the case of the two systems of this subsection: $\YC(v(u),z)=Y^{\mathrm C}_{\tilde A}$ by \cref{lem:gau-uniq}, both being solutions of $z\de_uY=\tilde AY$ with value $-\mathbf e_1$ (\cref{lem:gau-vanmod,lem:gau-vanq}~(c)). Put $\tilde Y:=Y^{\mathrm C}_{\tilde A}\in\bbQ[z^{-1}]\fps u^5$. The product $P\tilde Y$ lies in $\bbQ\lser z\fps u^5$: at each order in $u$ it is a finite sum of products of a power series in $z$ and a polynomial in $z^{-1}$. It solves the first system: $z\de_u(P\tilde Y)=(AP-P\tilde A)\tilde Y+P\tilde A\tilde Y=A(P\tilde Y)$. Its value at $u=0$ is $P(0,z)(-\mathbf e_1)=-P_{11}(0,z)\mathbf e_1$ by \cref{lem:gau-R1}. $P_{11}(0,z)Y^{\mathrm C}_A$ is also a solution, $P_{11}(0,z)$ being constant in $u$, with the same value by \cref{lem:gau-vanq}~(c). \Cref{lem:gau-uniq} with $B=A$ gives the identity.

(a) $P_{11}(0,z)$ is a unit of $\bbQ\fps{z^5}$ by \cref{prop:gau-graded}~(iii), so $\Pc$ is a graded regular gauge by \cref{prop:gau-graded}~(ii), and $(\Pc)_{11}(0,z)=1$; by \eqref{eq:gau-G}, $\Pc\YC(v(u))=\YCQ$. Conversely, let $P'$ be any regular gauge with $P'\YC(v(u))=\YCQ$. By \cref{thm:gau-U}~(b), (d), $P'=\Pc\cn(z)$ for some $\cn\in\bbQ\fps z^\times$, so that $\YCQ=\cn(z)\Pc\YC(v(u))=\cn(z)\YCQ$, i.e.\ $(\cn(z)-1)\YCQ=0$ in $\bbQ\lser z\fps u^5$; the coefficient of $u^0$ is $-(\cn(z)-1)\mathbf e_1$, so $\cn=1$. (b) If $P_{11}(0,z)=\mu$ is constant, then $P\YC(v(u))=\mu\YCQ$ lies in $\bbQ[z^{-1}]\fps u^5$. Conversely, the coefficient of $u^0$ of $P\YC(v(u))$ is $-P_{11}(0,z)\mathbf e_1$ with $P_{11}(0,z)\in\bbQ\fps z$, so it has no positive power of $z$ only if $P_{11}(0,z)$ is constant. With $P=\Pc\cn(z)$ we have $P_{11}(0,z)=\cn(z)$, which gives the last assertion.
\end{proof}

\begin{defn}[The canonical gauge]
\label{def:gau-can}
In a conifold frame, the \emph{canonical gauge} $\Pc$ is the regular gauge from the pulled-back model to the quintic system with $(\Pc)_{11}(0,z)=1$. It exists and is unique: $\Pc=P\,P_{11}(0,z)^{-1}$ for any graded regular gauge $P$ (\cref{thm:gau-G}~(a)), two graded regular gauges differ by a unit of $\bbQ\fps{z^5}$ (\cref{prop:gau-graded}), and a regular gauge $P'=\Pc\cn(z)$ (\cref{thm:gau-U}~(b)) with $P'_{11}(0,z)=1$ has $\cn=1$. Equivalently, $\Pc$ is the unique regular gauge with $\Pc\YC(v(u))=\YCQ$. It is graded, rational (\cref{thm:gau-U}~(d)), colour covariant, $\operatorname{Rot}_i\Pc=\Pc$ (\cref{prop:gau-graded}~(iv)), and $p_{0,a}(0)=1$ for its diagonal $z^0$-term (\cref{lem:gau-diag}). We write $\Pc=\sum_kz^k\sum_bp_{k,b}(u)E_{b+k,b}$.
\end{defn}

%-------------------------------------------------------------------------------
\subsection{Symplecticity}
\label{sec:gau-symp}

\begin{prop}[Symplecticity]
\label{prop:gau-symp}
Let $A,\tilde A$ be self-dual cyclic systems with one simple pole over $\Bbbk\fps u$ with (EP) and $r=\tilde r$, and let $P$ be a regular gauge from $\tilde A$ to $A$, $\cn(z):=P_{11}(0,z)$. Then
\begin{equation}
\label{eq:gau-symp}
P(u,-z)^T\eta\,P(u,z)=\cn(z)\cn(-z)\,\eta .
\end{equation}
In particular, for the two systems of \cref{sec:gau-van}, $\Pc$ is symplectic, $\Pc(u,-z)^T\eta\Pc(u,z)=\eta$; the symplectic regular gauges are the $\Pc\cn(z)$ with $\cn(z)\cn(-z)=1$, so symplecticity constrains only the even part of $\log(\cn/\cn(0))$; for graded gauges, $\cn=\pm\exp(a_1z^5+a_3z^{15}+a_5z^{25}+\cdots)$. Among the symplectic regular gauges, only (G) singles out $\Pc$.
\end{prop}

\begin{proof}
Put $S:=P(u,-z)^T\eta P(u,z)$. Replacing $z$ by $-z$ in the gauge equation gives $zP'(-z)=-AP(-z)+P(-z)\tilde A$, so $zP'(-z)^T=-P(-z)^TA^T+\tilde A^TP(-z)^T$. Hence
\[
zS'=\bigl(-P(-z)^TA^T+\tilde A^TP(-z)^T\bigr)\eta P(z)+P(-z)^T\eta\bigl(AP(z)-P(z)\tilde A\bigr)=\tilde A^TS-S\tilde A,
\]
because $A^T\eta=\eta A$ by self-duality (\cref{def:gau-onepole}). With $T:=\eta^{-1}S$ and $\tilde A^T\eta=\eta\tilde A$ this is $zT'=\tilde AT-T\tilde A$, and $T$ has coefficients in $\mathfrak{gl}_5(\Bbbk\fps u)$. By \cref{thm:gau-U}~(b) for the pair $(\tilde A,\tilde A)$, $T=t(z)\,1$ for some $t\in\Bbbk\fps z$; no invertibility of $T_0$ is needed. Compare the entries $(2,1)$ at $u=0$: by \cref{lem:gau-R1} column $1$ of $P(0,z)$ is $\cn(z)\mathbf e_1$, so $S_{21}(0,z)=\sum_aP_{a2}(0,-z)\eta_{a1}\cn(z)=5P_{22}(0,-z)\cn(z)=5\cn(-z)\cn(z)$, using $\eta_{a1}=5\delta_{a2}$ and $P_{22}(0,\cdot)=P_{11}(0,\cdot)$ ($r=\tilde r$); on the other hand $S_{21}=t(z)\eta_{21}=5t(z)$. So $t(z)=\cn(z)\cn(-z)$. For the consequences: $(\Pc)_{11}(0,z)=1$; every regular gauge is $\Pc\cn(z)$ with $(\Pc\cn)_{11}(0,z)=\cn(z)$; $\cn(z)\cn(-z)=1$ forces $\cn(0)=\pm1$ and is equivalent to the oddness of $\log(\cn/\cn(0))$; for graded gauges $\cn\in\Bbbk\fps{z^5}$.
\end{proof}

\begin{cor}[The dual identity]
\label{cor:gau-Gstar}
As row vectors in $\bbQ\lser z\fps u$,
\[
\YCQ(u,-z)^T\eta\,\Pc(u,z)=\YC(v(u),-z)^T\eta .
\]
More generally, $\YCQ(-z)^T\eta P(z)=P_{11}(0,z)\,\YC(-z)^T\eta$ for every regular gauge $P$.
\end{cor}

\begin{proof}
For a regular gauge, $P_{11}(0,z)$ is a unit of $\bbQ\fps z$, its constant term being $(P_0)_{11}(0)\neq0$. Substituting $-z$ for $z$ (a ring automorphism of $\bbQ\lser z\fps u$) in \eqref{eq:gau-G} gives $\YCQ(-z)=P_{11}(0,-z)^{-1}P(-z)\YC(-z)$. Transposing, multiplying on the right by $\eta P(z)$ and using \eqref{eq:gau-symp},
\[
\YCQ(-z)^T\eta P(z)=P_{11}(0,-z)^{-1}P_{11}(0,z)P_{11}(0,-z)\,\YC(-z)^T\eta=P_{11}(0,z)\,\YC(-z)^T\eta .\qedhere
\]
\end{proof}

% Section 3, part 6 of 7: HPF, polynomiality, formal WKB uniqueness, colour covariance of the model columns.

%-------------------------------------------------------------------------------
\subsection{The normalisation series}
\label{sec:gau-R0}

This subsection compares, colour by colour, the canonical gauge applied to the model columns with the rows of \cref{def:set-rows}. Colour $0$ (\cref{thm:gau-R0}) uses no (QRR); the other colours follow from colour $0$ by the colour covariance of the model columns, of the actual rows and of $\Pc$ (\cref{cor:gau-colours}). In this subsection $\theta_u:=u\,\dd/\dd u$ and $\theta_{\zh}:=\zh\,\dd/\dd\zh$, and a series $\Phi(\zh,u)=\sum_n\zh^n\mathrm{ph}_n(u)$ with $\mathrm{ph}_n\in\bbQ\lser s$ is identified with the series $\sum_nz^ns^{-4n}\mathrm{ph}_n$ in $z$, through $\zh=z/s^4=L^4z$ (\cref{def:set-points}). $\Dq$ acts on $\bbQ\lser s$ as in \cref{lem:gau-confexp}.

\begin{lem}[HPF]
\label{lem:gau-HPF}
For $e\in\bbZ$ let $\mathscr O_e:=1+\zh(1-u)\bigl(\tfrac45\theta_{\zh}-\tfrac e5-\theta_u\bigr)$, and for $\Phi=\sum_n\zh^n\mathrm{ph}_n$ as above put $\Phi^{[0]}:=\Phi$ and $\Phi^{[j]}:=\mathscr O_{-j}\Phi^{[j-1]}$, $j=1,\dots,5$.
\begin{enumerate}
\item $L\Phi$ satisfies $\PFfive^{(0)}(L\Phi)=0$, where $\PFfive^{(0)}$ is $\PFfive$ with $z\Dq$ replaced by $z\Dq+L$ (that is, $\re^{u_0/z}L\Phi$ solves $\PFfive$ formally), if and only if
\begin{equation}
\label{eq:gau-HPF}
\Phi^{[5]}-\Phi-(1-u)\Bigl[3\zh\Phi^{[4]}+\tfrac{17}5u\zh^2\Phi^{[3]}+\tfrac95u^2\zh^3\Phi^{[2]}+\tfrac{274}{625}u^3\zh^4\Phi^{[1]}+\tfrac{24}{625}u^4\zh^5\Phi\Bigr]=0 .
\end{equation}
\item The coefficient of $\zh^{n+1}$ of the left side of \eqref{eq:gau-HPF} is $(1-u)\bigl[(4n-5\theta_u)\mathrm{ph}_n+r_n\bigr]$, where, with $\Phi^\flat_{j,n}$ the coefficient of $\zh^n$ of $\Phi^{[j]}$ computed with $\mathrm{ph}_n$ replaced by $0$ (and $\Phi^{[j]}_m:=\coef{\zh^m}\Phi^{[j]}$, zero for $m<0$),
\[
r_n:=\sum_{j=1}^5\Bigl(\frac{4n+j}5-\theta_u\Bigr)\Phi^\flat_{j-1,n}-3\Phi^\flat_{4,n}-\tfrac{17}5u\Phi^{[3]}_{n-1}-\tfrac95u^2\Phi^{[2]}_{n-2}-\tfrac{274}{625}u^3\Phi^{[1]}_{n-3}-\tfrac{24}{625}u^4\mathrm{ph}_{n-4}
\]
depends only on $\mathrm{ph}_0,\dots,\mathrm{ph}_{n-1}$. So \eqref{eq:gau-HPF} holds if and only if $(4n-5\theta_u)\mathrm{ph}_n=-r_n$ for all $n\geq0$; here $r_0=0$, and $r_1=\frac35$ when $\mathrm{ph}_0=1$.
\item If every $\mathrm{ph}_n$ lies in $\bbQ\fps u$, \eqref{eq:gau-HPF} reads $(4n-5k)\coef{u^k}\mathrm{ph}_n=-\coef{u^k}r_n$ for all $n,k\geq0$: for $5k\neq4n$ the coefficient $\coef{u^k}\mathrm{ph}_n$ is fixed by $\mathrm{ph}_0,\dots,\mathrm{ph}_{n-1}$, while for $n=5m$ the coefficient $\fc_m:=\coef{u^{4m}}\mathrm{ph}_{5m}$ is free and $\coef{u^{4m}}r_{5m}=0$ is required. In particular $\mathrm{ph}_0$ is a constant, and $\mathrm{ph}_0=1$ forces $\mathrm{ph}_1=-\frac3{20}$.
\end{enumerate}
\end{lem}

\begin{proof}
(1) Let $F^\flat=\sum_m\zh^mF^\flat_m$ (a local letter) and $F:=s^eF^\flat=\sum_mz^ms^{e-4m}F^\flat_m$. Since $\Dq=-(1-u)u^{-1}\theta_u$ and $\theta_u(s^c)=\frac c5s^c$, $\Dq(s^cf)=(1-u)s^{c-5}\bigl(-\frac c5-\theta_u\bigr)f$. With $c=e-4m+4$, the coefficient of $z^m$ in $(z\Dq+L)F$ is $\Dq(s^{e-4m+4}F^\flat_{m-1})+s^{e-4m-1}F^\flat_m=s^{e-4m-1}\bigl[F^\flat_m+(1-u)\bigl(\frac{4(m-1)-e}5-\theta_u\bigr)F^\flat_{m-1}\bigr]$, which is the coefficient of $z^m$ of $s^{e-1}\mathscr O_eF^\flat$. So $(z\Dq+L)(s^eF^\flat)=s^{e-1}\mathscr O_eF^\flat$.
(2) By induction, $(z\Dq+L)^j(L\Phi)=(z\Dq+L)^j(s^{-1}\Phi)=s^{-1-j}\Phi^{[j]}$, and hence $z^m(z\Dq+L)^{5-m}(L\Phi)=s^{4m}\zh^m\,s^{m-6}\Phi^{[5-m]}=s^{-1}u^{m-1}\zh^m\Phi^{[5-m]}$.
(3) As $5^5q=1-u$ and $\prod_{k=1}^5(5M+kz)=5^5\sum_{j=0}^5\sigma_j\,z^jM^{5-j}$ for $M=z\Dq$, with $\sigma_j$ the elementary symmetric functions of $\frac15,\dots,\frac55$ (local letters), i.e.\ $\sigma_0=1$ and $(\sigma_1,\dots,\sigma_5)=(3,\frac{17}5,\frac95,\frac{274}{625},\frac{24}{625})$, we have $\PFfive=uM^5-1-(1-u)\bigl(3zM^4+\frac{17}5z^2M^3+\frac95z^3M^2+\frac{274}{625}z^4M+\frac{24}{625}z^5\bigr)$. Replacing $M$ by $z\Dq+L$, applying the result to $L\Phi$, using (2) and dividing by $s^{-1}$ gives the left side of \eqref{eq:gau-HPF}. This proves (i).
(4) From the definition, $\Phi^{[j]}_m=\Phi^{[j-1]}_m+(1-u)\bigl(\frac{4(m-1)+j}5-\theta_u\bigr)\Phi^{[j-1]}_{m-1}$: every raising of the $\zh$-degree carries a left factor $1-u$. Hence $\Phi^{[5]}_{n+1}-\mathrm{ph}_{n+1}=(1-u)\sum_{j=1}^5(\frac{4n+j}5-\theta_u)\Phi^{[j-1]}_n$. As $\Phi^{[j]}_n=\mathrm{ph}_n+\Phi^\flat_{j,n}$, the part of the coefficient of $\zh^{n+1}$ that contains $\mathrm{ph}_n$ is $(1-u)\bigl[\sum_{j=1}^5(\frac{4n+j}5-\theta_u)-3\bigr]\mathrm{ph}_n=(1-u)(4n-5\theta_u)\mathrm{ph}_n$, the $-3$ coming from $-3(1-u)\Phi^{[4]}_n$; the remaining terms are $(1-u)r_n$. Since $1-u$ is a unit, this proves (ii) apart from the values. For $n=0$ every term of $r_0$ vanishes. For $n=1$ and $\mathrm{ph}_0=1$: $\Phi^\flat_{j,1}=(1-u)\frac{j(j+1)}{10}$, and with $(c-\theta_u)(1-u)=c(1-u)+u$,
\[
r_1=\tfrac15\bigl(\tfrac65(1-u)+u\bigr)+\tfrac35\bigl(\tfrac75(1-u)+u\bigr)+\tfrac65\bigl(\tfrac85(1-u)+u\bigr)+2\bigl(\tfrac95(1-u)+u\bigr)-6(1-u)-\tfrac{17}5u=\tfrac35 .
\]
(5) For power series, $\theta_u$ acts on $u^k$ by $k$, which gives the first part of (iii); the eigenvalue $4n-5k$ vanishes only for $n=5m$, $k=4m$. For $n=0$, $(4\cdot0-5k)\coef{u^k}\mathrm{ph}_0=0$ gives $\mathrm{ph}_0\in\bbQ$; for $n=1$, $(4-5k)\coef{u^k}\mathrm{ph}_1=-\frac35\delta_{k0}$ gives $\mathrm{ph}_1=-\frac3{20}$.
\end{proof}

\begin{lem}[Power-series solutions of HPF are polynomial]
\label{lem:gau-poly}
Let $\Phi$ be a solution of \eqref{eq:gau-HPF} with $\mathrm{ph}_0=1$ and every $\mathrm{ph}_n\in\bbQ\fps u$.
\begin{enumerate}
\item Every $\mathrm{ph}_n$ is a polynomial in $u$.
\item For $\mathsf k\in\bbQ\fps z$ with $\mathsf k(0)=1$, the product $\Phi\,\mathsf k(z)$ (with $z^n=s^{4n}\zh^n$) solves \eqref{eq:gau-HPF}, and its coefficients lie in $\bbQ\fps u$ if and only if $\mathsf k\in\bbQ\fps{z^5}$.
\item The solutions of \eqref{eq:gau-HPF} with $\mathrm{ph}_0=1$ and coefficients in $\bbQ\fps u$ are exactly the $\Phi\,\mathsf k(z)$, $\mathsf k\in1+z^5\bbQ\fps{z^5}$; for every choice of the free coefficients $(\fc_m)_{m\geq1}$ there is exactly one of them, and $\Phi^\circ$ denotes the one with all $\fc_m=0$.
\end{enumerate}
In particular, as soon as one solution with coefficients in $\bbQ\fps u$ and $\mathrm{ph}_0=1$ exists, the resonance conditions $\coef{u^{4m}}r_{5m}=0$ hold for every choice of the free coefficients.
\end{lem}

\begin{proof}
(i) By induction on $n$: if $\mathrm{ph}_0,\dots,\mathrm{ph}_{n-1}$ are polynomials, so is $r_n$, which is obtained from them by $\theta_u$ and multiplication by polynomials in $u$; by \cref{lem:gau-HPF}~(iii), $\coef{u^k}\mathrm{ph}_n=-\coef{u^k}r_n/(4n-5k)$ vanishes for $k>\deg r_n$, except possibly for the single free coefficient at $k=4n/5$. (ii) In the variables $(\zh,u)$, $z^m=s^{4m}\zh^m$ satisfies $\theta_{\zh}z^m=mz^m$ and $\theta_uz^m=\frac{4m}5z^m$, so $(\frac45\theta_{\zh}-\theta_u)(z^mF)=z^m(\frac45\theta_{\zh}-\theta_u)F$; hence every $\mathscr O_e$ commutes with multiplication by a series in $z$, and so does the left side of \eqref{eq:gau-HPF}, whose other coefficients involve only $\zh$ and $u$. The coefficient of $\zh^N$ of $\Phi\mathsf k$ is $\sum_n\mathsf k_ns^{4n}\mathrm{ph}_{N-n}$; if $n_0$ is the least index with $5\nmid n_0$ and $\mathsf k_{n_0}\neq0$, the coefficient of $\zh^{n_0}$ contains $\mathsf k_{n_0}s^{4n_0}$, a fractional power of $u$, plus power series; so the coefficients lie in $\bbQ\fps u$ if and only if no such $n_0$ exists. (iii) Let $\Phi'$ be a second such solution. We construct $\mathsf k^{(N)}\in1+z^5\bbQ[z^5]$ with $\Phi'\equiv\Phi\mathsf k^{(N)}$ modulo $\zh^N$, starting from $\mathsf k^{(1)}=1$. Both $\Phi'$ and $\Phi\mathsf k^{(N)}$ are solutions with coefficients in $\bbQ\fps u$ and $\mathrm{ph}_0=1$, by (ii), and they agree below $\zh^N$, so they have the same $r_N$, and the difference $\delta$ of their coefficients of $\zh^N$ satisfies $(4N-5\theta_u)\delta=0$. If $5\nmid N$, then $\delta=0$ and we put $\mathsf k^{(N+1)}:=\mathsf k^{(N)}$. If $N=5m$, then $\delta=\delta_0u^{4m}$ with $\delta_0\in\bbQ$, and $\mathsf k^{(N+1)}:=\mathsf k^{(N)}+\delta_0z^{5m}$ works, since $z^{5m}\Phi=u^{4m}\zh^{5m}(1+O(\zh))$. The limit $\mathsf k$ satisfies $\Phi'=\Phi\mathsf k$; it is unique because $\Phi$ is a unit of $\bbQ\lser s\fps{\zh}$. Conversely $\Phi\mathsf k$ is such a solution by (ii). The free coefficients of $\Phi\mathsf k$ are $\fc_m(\Phi\mathsf k)=\mathsf k_{5m}+\sum_{0<m'<m}\mathsf k_{5m'}\fc_{m-m'}(\Phi)+\fc_m(\Phi)$, a triangular system in the $\mathsf k_{5m}$ with unit diagonal; so every sequence of free coefficients is attained exactly once. The last assertion follows: for every choice of the free coefficients a solution exists, so no resonance condition fails.
\end{proof}

\begin{lem}[Formal WKB uniqueness]
\label{lem:gau-wkb}
Let $\Bbbk\supseteq\bbQ(\zeta)$ be a field (for $i=0$, any field $\Bbbk\supseteq\bbQ$), and let $(\cR,\Dq,L)$ be $(\Bbbk\fps q,\,q\,\dd/\dd q,\,(1-5^5q)^{-1/5})$ or $(\Bbbk\lser s,\Dq,s^{-1})$; in both, $\Dq L=(L^6-L)/5$ and $5^5q=1-L^{-5}$. For $i\in\bbZ/5$ let $\PFfive^{(i)}\in\cR[z][\Dq]$ be $\PFfive$ with $z\Dq$ replaced by $z\Dq+\zeta^iL$; a \emph{colour-$i$ WKB amplitude} of $\PFfive$ over $\cR$ is a series $\cA=\sum_{n\geq0}\cA_nz^n\in\cR\fps z$ with $\PFfive^{(i)}\cA=0$ (so that $\re^{u_i/z}\cA$ solves $\PFfive$ formally, $\Dq u_i=\zeta^iL$).
\begin{enumerate}
\item $\PFfive^{(i)}=\sum_{k=1}^5z^k\mathcal T^{\mathrm w}_k$ with $\mathcal T^{\mathrm w}_k\in\cR[\Dq]$ (local letters), and $\mathcal T^{\mathrm w}_1(Lg)=5\zeta^{4i}\Dq g$ for $g\in\cR$.
\item If $\cA,\cA'$ are colour-$i$ WKB amplitudes with $\cA_0=\cA'_0=L$, then $\cA'=\cA\,\mathsf k(z)$ for a unique $\mathsf k\in1+z\Bbbk\fps z$; for $i=0$ and $\cA,\cA'$ with coefficients in $\bbQ\fps q$ or $\bbQ\lser s$, $\mathsf k\in1+z\bbQ\fps z$.
\item $\cA(z)\mapsto\cA(\zeta^{-i}z)$ is a bijection from colour-$0$ to colour-$i$ WKB amplitudes that preserves $\cA_0$.
\end{enumerate}
\end{lem}

\begin{proof}
(i) The coefficient of $z^0$ of $\PFfive^{(i)}$ is the multiplication by $(\zeta^iL)^5-1-q(5\zeta^iL)^5=L^5(1-5^5q)-1=0$. The coefficient of $z$ in $(z\Dq+\zeta^iL)^5$ is $\sum_{j=0}^4(\zeta^iL)^j\Dq(\zeta^iL)^{4-j}$, which acts as
\[
f\mapsto\zeta^{4i}\bigl(5L^4\Dq f+10L^3(\Dq L)f\bigr)=\zeta^{4i}\bigl(5L^4\Dq f+2(L^9-L^4)f\bigr).
\]
The coefficient of $z$ in $\prod_{k=1}^5\bigl(5\zeta^iL+z(5\Dq+k)\bigr)$ is $\sum_k(5\zeta^iL)^{k-1}(5\Dq+k)(5\zeta^iL)^{5-k}$, which acts as $f\mapsto5^4\zeta^{4i}(25L^4\Dq f+10L^9f+5L^4f)$. With $5^4q=(1-L^{-5})/5$ the difference is $\mathcal T^{\mathrm w}_1f=\zeta^{4i}\bigl(5L^{-1}\Dq f+(L^{-1}-L^4)f\bigr)$, and $\mathcal T^{\mathrm w}_1(Lg)=\zeta^{4i}\bigl(5\Dq g+(L^5-1)g+(1-L^5)g\bigr)=5\zeta^{4i}\Dq g$.
(ii) The coefficient of $z^{n+1}$ of $\PFfive^{(i)}\cA$ is $\sum_{k=1}^5\mathcal T^{\mathrm w}_k\cA_{n+1-k}$, so the equation is $\mathcal T^{\mathrm w}_1\cA_n=-\sum_{k=2}^5\mathcal T^{\mathrm w}_k\cA_{n+1-k}$. As $L$ is a unit of $\cR$ and the constants of $\Dq$ on $\cR$ are $\Bbbk$, (i) gives $\ker\mathcal T^{\mathrm w}_1=\Bbbk L$. Multiplication by a constant series $\mathsf k(z)$ commutes with $\PFfive^{(i)}$. Suppose $\cA'\equiv\cA\mathsf k^{(n)}$ modulo $z^n$ for a polynomial $\mathsf k^{(n)}$ of degree less than $n$ with $\mathsf k^{(n)}(0)=1$ (for $n=1$, $\mathsf k^{(1)}=1$). Then $B:=\cA'-\cA\mathsf k^{(n)}$ is annihilated by $\PFfive^{(i)}$ and $B\equiv0$ modulo $z^n$, so its equation at $z^{n+1}$ is $\mathcal T^{\mathrm w}_1B_n=0$, i.e.\ $B_n=\mathsf k_nL$ with $\mathsf k_n\in\Bbbk$; and $\mathsf k^{(n+1)}:=\mathsf k^{(n)}+\mathsf k_nz^n$ gives $\cA'\equiv\cA\mathsf k^{(n+1)}$ modulo $z^{n+1}$, since $\coef{z^n}(\mathsf k_nz^n\cA)=\mathsf k_nL$. The limit $\mathsf k$ satisfies $\cA'=\cA\mathsf k$, and it is unique because $\cA$ is a unit. In the rational case of colour $0$ all data are rational, so every $\mathsf k_n$ is.
(iii) Let $\sigma_i(\cA)(z):=\cA(\zeta^{-i}z)$, a ring automorphism of $\cR\fps z$ commuting with $\Dq$. Then $\sigma_i\circ z=\zeta^{-i}z\circ\sigma_i$ and $\sigma_i\circ(z\Dq+L)=\zeta^{-i}(z\Dq+\zeta^iL)\circ\sigma_i$. Apart from its constant term $-1$, $\PFfive=M^5-1-q\prod_k(5M+kz)$ is homogeneous of degree $5$ in $(M,z)$, $M=z\Dq$; as $\zeta^{-5i}=1$, $\sigma_i\circ\PFfive^{(0)}=\PFfive^{(i)}\circ\sigma_i$. So $\sigma_i$ maps colour-$0$ amplitudes to colour-$i$ amplitudes, with inverse $\sigma_{-i}$, and it fixes $\cA_0$.
\end{proof}

By (iii), (ii) for colour $i$ also follows from (ii) for colour $0$.

\begin{lem}[Colour covariance of the model columns]
\label{lem:gau-colmod}
For $i\in\bbZ/5$ let $F^{(i)}=(F^{(i)}_b)_{b\in\bbZ/5}$ be the solution with $F^{(i)}_0(0)=1$ of the colour-$i$ T-system of \cref{prop:set-modrows}~(b), with $\mathcal T^{(i)}_k=-\zeta^i+\frac t5(k-4t\,\dd/\dd t)$; thus $\Fs_b=F^{(0)}_b$.
\begin{enumerate}
\item $F^{(i)}_b(t)=\zeta^{ib}\Fs_b(\zeta^{-i}t)$.
\item The rows $\Rhat^{\mathrm m,(i)}_b(t)=\zeta^{i(b+1)}\Fs_b(\zeta^{-i}t)=\zeta^iF^{(i)}_b(t)$ of \cref{def:set-M0} are the rows of a WKB solution of the model system with exponential $\re^{U_i/z}$, $U_i:=-\frac54\zeta^i\varsigma^4$: the \emph{model columns} $\Qr^{\mathrm m,(i)}_b(v,z):=\varsigma^{\nu_b-1}\Rhat^{\mathrm m,(i)}_b(z/\varsigma^4)$ satisfy $z\,\de_v\Qr^{\mathrm m,(i)}_b+(\de_vU_i)\Qr^{\mathrm m,(i)}_b=b^{\mathrm m}_b\Qr^{\mathrm m,(i)}_{b+1}$.
\item At $v=v(u)$: $\Qr^{\mathrm m,(i)}_b(u,z)=\zeta^{i(b+1)}\Qr^{\mathrm m,(0)}_b(u,\zeta^{-i}z)$, i.e.\ $\Qr^{\mathrm m,(i)}=\operatorname{Rot}_i\Qr^{\mathrm m,(0)}$.
\end{enumerate}
\end{lem}

\begin{proof}
(i) $t\,\dd/\dd t$ commutes with the substitution $t\mapsto\zeta^{-i}t$, and $t=\zeta^i(\zeta^{-i}t)$; hence $\mathcal T^{(i)}_k[G(\zeta^{-i}t)]=\zeta^i[\mathcal T^{(0)}_kG](\zeta^{-i}t)$ for every series $G$. Put $G_b(t):=\zeta^{ib}\Fs_b(\zeta^{-i}t)$. Then $G_1=\zeta^i[\mathcal T^{(0)}_{-1}\Fs_0](\zeta^{-i}t)=\mathcal T^{(i)}_{-1}G_0$, $G_2=-\zeta^{2i}[\mathcal T^{(0)}_{-2}\Fs_1](\zeta^{-i}t)=-\mathcal T^{(i)}_{-2}G_1$, and in the same way $G_3=\mathcal T^{(i)}_2G_2$ and $G_4=-\mathcal T^{(i)}_1G_3$; finally $G_0=-[\mathcal T^{(0)}_0\Fs_4](\zeta^{-i}t)=-\zeta^{-i}\mathcal T^{(i)}_0[\Fs_4(\zeta^{-i}t)]=-\zeta^{-5i}\mathcal T^{(i)}_0G_4=-\mathcal T^{(i)}_0G_4$, and $G_0(0)=1$. By the uniqueness in \cref{prop:set-modrows}~(b) (the pivot of the recursion at $\coef{t^{n+1}}$ is $4n\zeta^{4i}\neq0$), $F^{(i)}=G$. (ii) By \cref{prop:set-modrows}~(b), $\re^{U_i/z}\varsigma^{\nu_b-1}F^{(i)}_b(z/\varsigma^4)$ solves the model system, and so does its multiple by the constant $\zeta^i$. (iii) As $\varsigma$ depends on $u$ only, $\Qr^{\mathrm m,(i)}_b=\varsigma^{\nu_b-1}\zeta^{i(b+1)}\Fs_b(\zeta^{-i}z/\varsigma^4)=\zeta^{i(b+1)}\Qr^{\mathrm m,(0)}_b(u,\zeta^{-i}z)$.
\end{proof}

% Section 3, part 7 of 7: colour covariance of the actual rows, the normalisation series D, the rows in the ring, the colours of D, (F_g).

\begin{lem}[Colour covariance of the actual rows]
\label{lem:gau-colrows}
For every $i\in\bbZ/5$, as identities of formal series in $z$ with coefficients in $\bbQ(\zeta)\fps q$,
\[
\Phi_i(\zh,u)=\Phi_0(\zeta^{-i}\zh,u)\qquad\text{and}\qquad \Qr^{(i)}_a(z)=\zeta^{i(a+1)}\Qr^{(0)}_a(\zeta^{-i}z),\ \text{i.e.}\ \Qr^{(i)}_{a,n}=\zeta^{i(a+1-n)}\Qr^{(0)}_{a,n},
\]
where $\Qr^{(i)}_{a,n}:=\coef{z^n}\Qr^{(i)}_a$ are the coefficients of the rows of \cref{def:set-rows}.
\end{lem}

\begin{proof}
The proof uses the forward relations of \cref{prop:set-cyclicLR} (only their shape: the coefficients $C_1\kappa^{\mathrm{LR}}_a$ do not depend on $z$ or $i$), $\Dq u_i=\zeta^iL$ and the WKB form of row $0$ (\cref{cor:set-cyclicPF5}), (QRR) in the form of \cref{thm:set-qrr}, \cref{lem:gau-wkb} and \cref{lem:set-hrot}. All objects are formal series in $z$ over $\bbQ(\zeta)\fps q$.

\emph{Step 1.} By \cref{cor:set-cyclicPF5}~(b), $\re^{u_i/z}zL\Phi_i$ solves $\PFfive$; as $z$ is a constant, $L\Phi_i$ is a colour-$i$ WKB amplitude over $\bbQ(\zeta)\fps q$ in the sense of \cref{lem:gau-wkb}. Its coefficient of $z^0$ is $L$, because $R=1+O(z)$ gives $\coef{z^0}\Qr^{(i)}_0=\eta(1,\tilde e_i)=\zeta^iL/C_0$ (\cref{cor:set-cyclicPF5}~(a), \cref{def:set-frob}). By \cref{lem:gau-wkb}~(iii), the colour-$0$ amplitude with $z$ replaced by $\zeta^{-i}z$, which is $L\Phi_0(\zeta^{-i}\zh,u)$ since $\zh=L^4z$, is a colour-$i$ WKB amplitude with coefficient $L$ at $z^0$.

\emph{Step 2.} By \cref{lem:gau-wkb}~(ii), $L\Phi_0(\zeta^{-i}\zh,u)=L\Phi_i(\zh,u)\,\mathsf k(z)$ for some $\mathsf k\in1+z\bbQ(\zeta)\fps z$.

\emph{Step 3.} Set $q=0$, where $u=1$, $L=1$, $\zh=z$ and $C_0=1$; as $\mathsf k$ does not depend on $q$, $\Phi_0(\zeta^{-i}z,1)=\Phi_i(z,1)\,\mathsf k(z)$. By (QRR) (\cref{thm:set-qrr}), $\Phi_i(z,1)=h_{\bw_i}(z)$ and $\Phi_0(\zeta^{-i}z,1)=h_{\bw_0}(\zeta^{-i}z)$, which is $h_{\bw_i}(z)$ by \cref{lem:set-hrot}. As $h_{\bw_i}\in1+z\bbQ(\zeta)\fps z$ is a unit, $\mathsf k=1$; this proves the first identity. This step is the only place in this section where (QRR) enters.

\emph{Step 4.} Put $X^{(i)}_a:=\zeta^{-i}\Qr^{(i)}_a$, so $X^{(i)}_0=L\Phi_i/C_0$ by the definition of $\Phi_i$ (\cref{def:set-rows}), and $X^{(i)}_0(z)=X^{(0)}_0(\zeta^{-i}z)$ by Step 3. The forward relations of \cref{prop:set-cyclicLR} are $X^{(i)}_{a+1}=(C_1\kappa^{\mathrm{LR}}_a)^{-1}(z\Dq+\zeta^iL)X^{(i)}_a$, $a=0,\dots,3$. For a series $F$, $(z\Dq+\zeta^iL)[F(\zeta^{-i}z)]=\zeta^i\bigl[(z\Dq+L)F\bigr](\zeta^{-i}z)$. So if $X^{(i)}_a(z)=\zeta^{ia}X^{(0)}_a(\zeta^{-i}z)$, then $X^{(i)}_{a+1}(z)=\zeta^{ia}(C_1\kappa^{\mathrm{LR}}_a)^{-1}\zeta^i\bigl[(z\Dq+L)X^{(0)}_a\bigr](\zeta^{-i}z)=\zeta^{i(a+1)}X^{(0)}_{a+1}(\zeta^{-i}z)$, because $C_1\kappa^{\mathrm{LR}}_a$ depends neither on $z$ nor on $i$.

\emph{Step 5.} Hence $\Qr^{(i)}_a(z)=\zeta^iX^{(i)}_a(z)=\zeta^{i(a+1)}X^{(0)}_a(\zeta^{-i}z)=\zeta^{i(a+1)}\Qr^{(0)}_a(\zeta^{-i}z)$, and taking the coefficient of $z^n$ gives the phase $\zeta^{i(a+1-n)}$.
\end{proof}

This is the colour covariance of the rows announced after \cref{def:set-rows}. That it survives the conifold expansion and hatting is part of \cref{cor:gau-rowsNN}.

\begin{thm}[The normalisation series $\Dn$; colour $0$]
\label{thm:gau-R0}
Let $(\alpha,\beta)$ be a conifold frame, $\Pc$ the canonical gauge, and $\Qr^{\mathrm m,(i)}$ the model columns of \cref{lem:gau-colmod} at $v=v(u)$.
\begin{enumerate}
\item For every $i$, $\Pc\Qr^{\mathrm m,(i)}$ is a formal WKB solution of the quintic system with exponential $\re^{u_i/z}$: its rows $X_a\in\bbQ(\zeta)\lser s\fps z$ satisfy $(z\Dq+\zeta^iL)X_a=\kappa_aX_{a+1}$ for $a\in\bbZ/5$, with $L=s^{-1}$. With $\Lambda_s:=\diag(s^{1-\nu_a})$, $\Lambda_\varsigma:=\diag(\varsigma^{1-\nu_b})$ and $\Ghat(u,\zh):=\Lambda_s\Pc(u,s^4\zh)\Lambda_\varsigma^{-1}$ (the hatted gauge of \cref{lem:mai-Ghat}),
\[
\Lambda_s\,\Pc\,\Qr^{\mathrm m}\big|_{z=s^4\zh}=\Ghat(u,\zh)\,\Rhat^{\mathrm m}(\tau(u)\zh),\qquad \coef{\zh^k}\Ghat_{b+k,b}=\tau^{(1-\nu_b)/4}u^{\mu(b,k)}p_{k,b}(u),
\]
with $\mu(b,k)=(4k+\nu_b-\nu_{b+k})/5\in\bbZ$, $\mu(b,k)\geq\max(0,(4k-4)/5)$; this is a power series in $u$ whose coefficients are power series in $\zh$.
\item $\Phican(\zh,u):=\varpi_0\cdot\bigl[\Ghat(u,\zh)\Rhat^{\mathrm m}(\tau\zh)\bigr]_{00}$ (row $0$, colour $0$) solves \eqref{eq:gau-HPF}, with coefficients in $\bbQ\fps u$ and $\mathrm{ph}^{\mathrm{can}}_0=1$. Hence every $\mathrm{ph}^{\mathrm{can}}_n$ is a polynomial in $u$ (\cref{lem:gau-poly}) and $A^{\mathrm{can}}_n:=L^{4n+1}\mathrm{ph}^{\mathrm{can}}_n(L^{-5})$ lies in $\bbQ[L,L^{-1}]$.
\item Let $A^{\mathrm{act}}:=C_0\Qr^{(0)}_0=\sum_nA^{\mathrm{act}}_nz^n$ be the flat amplitude of the actual colour-$0$ row $0$ (\cref{def:set-rows}), with $A^{\mathrm{act}}_n\in\bbQ(\zeta)\fps q$. In $\bbQ(\zeta)\fps q$, with $L=(1-5^5q)^{-1/5}$, both $A^{\mathrm{can}}:=\sum_nA^{\mathrm{can}}_nz^n$ and $A^{\mathrm{act}}$ are colour-$0$ WKB amplitudes of $\PFfive$ with coefficient $L$ at $z^0$. Hence $A^{\mathrm{can}}=A^{\mathrm{act}}\,\Dn_0(z)$ for a unique $\Dn_0\in1+z\bbQ(\zeta)\fps z$, and $A^{\mathrm{act}}_n\in\bbQ(\zeta)[L,L^{-1}]$ for every $n$.
\item In $\bbQ(\zeta)[L,L^{-1}]\fps z$, embedded in $\bbQ(\zeta)\lser s\fps z$ by $L\mapsto s^{-1}$, the identity of (iii) holds; in hatted form it reads $\Phiact(\zh,u)\,\Dn_0(s^4\zh)=\Phican(\zh,u)$, $\Phiact:=\Phi_0=L^{-1}A^{\mathrm{act}}$. The rows of the actual colour-$0$ column have conifold expansions (\cref{cor:gau-rowsNN}), they satisfy $(z\Dq+L)X_a=\kappa_aX_{a+1}$, and
\begin{equation}
\label{eq:gau-R0zero}
\Qr^{(0)}\,\Dn_0(z)=\Pc\,\Qr^{\mathrm m,(0)}(v(u)),
\end{equation}
with $\Dn_0$ constant in $u$ and the same exponential $\re^{u_0/z}$ stripped on both sides.
\end{enumerate}
\end{thm}

\begin{proof}
(i) By \cref{lem:gau-colmod}~(ii), the amplitude $\Qr^{\mathrm m,(i)}$ of the exponential $\re^{U_i/z}$ solves the model system; at $v=v(u)$ the chain rule gives $z\de_u\Qr^{\mathrm m,(i)}+(\de_uU_i)\Qr^{\mathrm m,(i)}=\tilde A\Qr^{\mathrm m,(i)}$ with $U_i=U_i(v(u))$. By \cref{def:set-tau} and \cref{lem:set-canon}, $U_i(v(u))=-\frac54\zeta^iv^{4/5}=-\frac54\zeta^iu^{4/5}/\tau=u_i-u_i^{\mathrm c}$, so $\de_uU_i=\de_uu_i$. For $X:=\Pc\Qr^{\mathrm m,(i)}$ the gauge equation gives $z\de_uX+(\de_uu_i)X=(A\Pc-\Pc\tilde A)\Qr^{\mathrm m,(i)}+\Pc\tilde A\Qr^{\mathrm m,(i)}=AX$; multiplying by $-(1-u)$ and using $b_a=-\kappa_a/(1-u)$, $\Dq u_i=\zeta^iL$, this is $(z\Dq+\zeta^iL)X_a=\kappa_aX_{a+1}$. The amplitude of $\re^{u_i/z}$ is the same as that of $\re^{U_i/z}$, because $u_i-U_i=u_i^{\mathrm c}$ is a constant; the factor $\re^{u_i^{\mathrm c}/z}$, which is not a series in $z$, is not part of any amplitude. For the hatted form: $\varsigma=s\tau^{-1/4}$ (\cref{def:set-tau}), so $z/\varsigma^4=\tau\zh$ and $\Lambda_\varsigma\Qr^{\mathrm m}=\Rhat^{\mathrm m}(\tau\zh)$; and the entry $(b+k,b)$ of $\Lambda_s\Pc(u,s^4\zh)\Lambda_\varsigma^{-1}$ at $\zh^k$ is $s^{1-\nu_{b+k}}s^{4k}p_{k,b}\,s^{\nu_b-1}\tau^{(1-\nu_b)/4}=u^{\mu(b,k)}\tau^{(1-\nu_b)/4}p_{k,b}$. Since $\nu_a\equiv-a\pmod5$, $5\mid4k+\nu_b-\nu_{b+k}$; since $\nu$ takes values in $\{-1,\dots,3\}$, $\mu(b,k)\geq(4k-4)/5$, and $\mu(b,0)=0$. So $\Ghat$ has entries in $\bbQ[\zh]\fps u$, and $\Rhat^{\mathrm m}(\tau\zh)$ has entries in $\bbQ(\zeta)\fps u\fps{\zh}$.

(ii) As $\nu_0=0$, row $0$ of $\Lambda_s\Pc\Qr^{\mathrm m,(0)}$ is $s\,(\Pc\Qr^{\mathrm m,(0)})_0$, so $L\Phican=\varpi_0(\Pc\Qr^{\mathrm m,(0)})_0$. By (i), the rows of $\Pc\Qr^{\mathrm m,(0)}$ satisfy the cyclic system in the conjugated form of \cref{cor:gau-cycsys}~(ii) ($\cR=\bbQ\lser s\lser z$, $w=L$, $(y_0,y_1)=(\varpi_0,\varpi_1)$); hence $\varpi_0(\Pc\Qr^{\mathrm m,(0)})_0=L\Phican$ solves $\PFfive^{(0)}$, and $\Phican$ solves \eqref{eq:gau-HPF} by \cref{lem:gau-HPF}~(i). For colour $0$ all data are rational ($\Rhat^{\mathrm m,(0)}_b=\Fs_b$, \cref{def:set-M0}; $\Pc$ rational, \cref{def:gau-can}), so by (i) the coefficients of $\Phican$ lie in $\bbQ\fps u$. As $\Ghat(u,0)$ is diagonal with $(0,0)$ entry $\tau^{1/4}p_{0,0}=1/\varpi_0$ (\cref{lem:gau-diag}) and $\Fs_0(0)=1$, $\mathrm{ph}^{\mathrm{can}}_0=1$. Then \cref{lem:gau-poly}~(i) applies, and $A^{\mathrm{can}}_n=u^{-(4n+1)/5}\mathrm{ph}^{\mathrm{can}}_n(u)$ is a Laurent polynomial in $L=u^{-1/5}$.

(iii) \emph{The two rings.} $\Dq$ preserves $\bbQ(\zeta)[L,L^{-1}]$ ($\Dq L=(L^6-L)/5$, and $q=(1-L^{-5})/5^5$), and both $L\mapsto(1-5^5q)^{-1/5}\in\bbQ(\zeta)\fps q$ and $L\mapsto s^{-1}\in\bbQ(\zeta)\lser s$ are injective homomorphisms of differential rings: the second clearly, the first because $L-1\in5^4q+q^2\bbQ\fps q$ has order $1$ in $q$, so a non-zero polynomial in $L$ has a non-zero image, and $L$ is a unit. The colour-$0$ WKB equations of $A^{\mathrm{can}}$ are identities in $\bbQ[L,L^{-1}]$, coefficient by coefficient in $z$; they hold after $L\mapsto s^{-1}$ by (ii) (with $\zh=L^4z$, $\sum_nA^{\mathrm{can}}_nz^n=L\Phican$), hence in $\bbQ[L,L^{-1}]$, hence in $\bbQ\fps q$. Its coefficient of $z^0$ is $L\,\mathrm{ph}^{\mathrm{can}}_0=L$. $A^{\mathrm{act}}$ is a colour-$0$ WKB amplitude over $\bbQ(\zeta)\fps q$ by \cref{cor:set-cyclicPF5}~(b), with coefficient $C_0\eta(1,\tilde e_0)=L$ at $z^0$ (\cref{def:set-frob}). \Cref{lem:gau-wkb}~(ii), with $\Bbbk=\bbQ(\zeta)$, gives $A^{\mathrm{can}}=A^{\mathrm{act}}\Dn_0(z)$ with $\Dn_0\in1+z\bbQ(\zeta)\fps z$. Then $A^{\mathrm{act}}_n=\sum_{m\leq n}A^{\mathrm{can}}_m\coef{z^{n-m}}\Dn_0^{-1}$ is the image of a unique Laurent polynomial in $L$ over $\bbQ(\zeta)$.

(iv) By the injectivity in (iii) the identity $A^{\mathrm{can}}=A^{\mathrm{act}}\Dn_0$ holds in $\bbQ(\zeta)[L,L^{-1}]\fps z$, hence in $\bbQ(\zeta)\lser s\fps z$; dividing by $L$ and writing $z=s^4\zh$ gives the hatted form. By (iii) the rows of the actual colour-$0$ column have $\logq$-free representatives in $\Nring\otimes\bbQ(\zeta)$: row $0$ is $A^{\mathrm{act}}_n(L)c_0^{-1}$, and the rows $a\geq1$ follow by the forward relations of \cref{prop:set-cyclicLR}, whose coefficients are represented in $\Nring$ (\cref{prop:gau-quintic}~(i)); this is \cref{cor:gau-rowsNN}~(i) for colour $0$. By \cref{prop:gau-quintic}~(i) their conifold expansions satisfy $(z\Dq+L)X_a=\kappa_aX_{a+1}$ for $a=0,\dots,3$. Row $0$ of \eqref{eq:gau-R0zero}: the expansion of $\Qr^{(0)}_0$ is $A^{\mathrm{act}}/\varpi_0$ (with $L=s^{-1}$), and $A^{\mathrm{act}}\Dn_0/\varpi_0=A^{\mathrm{can}}/\varpi_0=L\Phican/\varpi_0=(\Pc\Qr^{\mathrm m,(0)})_0$ by (ii). The rows $a=1,\dots,4$ of both sides are obtained from row $0$ by the same relations $X_{a+1}=\kappa_a^{-1}(z\Dq+L)X_a$ (on the left because $\Dn_0$ is constant, on the right by (i)), so they agree. The closing relation $a=4$ of the left side follows from that of the right side.
\end{proof}

No normalisation of the actual rows is used in the theorem, and colour $0$ uses no (QRR): only the facts that $A^{\mathrm{act}}$ is a WKB amplitude with leading term $L$ and that the forward relations hold. For the same reason the theorem gives $\Dn_0$ only over $\bbQ(\zeta)$: a WKB amplitude is fixed by its leading term only up to a constant series (\cref{lem:gau-wkb}), and nothing used so far fixes that series to be rational. That $\Dn_0$, equivalently every $A^{\mathrm{act}}_n$, has rational coefficients follows from the value of $A^{\mathrm{act}}$ at $q=0$ given by (QRR); this is part of the glue \cref{thm:con-glue}, and it is not used before.

\begin{cor}[The rows lie in the differential ring]
\label{cor:gau-rowsNN}
\begin{enumerate}
\item For all $i,a,n$ the row coefficient $\Qr^{(i)}_{a,n}$ is the evaluation (\cref{def:set-confexp}) of an explicit $\logq$-free element of $\Nring\otimes\bbQ(\zeta)$: for colour $0$, row $0$ is $A^{\mathrm{act}}_n(L)\,c_0^{-1}$ with $A^{\mathrm{act}}_n\in\bbQ(\zeta)[L,L^{-1}]$ (\cref{thm:gau-R0}~(iii)); the rows $a=1,\dots,4$ follow from $\Qr_{a+1,n}=(C_1\kappa^{\mathrm{LR}}_a)^{-1}(\Dq\Qr_{a,n-1}+L\Qr_{a,n})$, with $(C_1\kappa^{\mathrm{LR}}_a)^{-1}$ represented by $(c_0^2/\mathsf W,\ \mathsf W^2L^{-5}/c_0^2,\ c_0^2/\mathsf W,\ c_0^{-1})$; and colour $i$ is $\zeta^{i(a+1-n)}$ times colour $0$ (\cref{lem:gau-colrows}). These representatives are used throughout; in this sense the rows lie in $\Nring\otimes\bbQ(\zeta)$ and have conifold expansions.
\item Consequently every polynomial over $\bbQ(\zeta)$ in finitely many row coefficients and in the $\Delta_i^{\pm1/2}=(\zeta^{4i}C_0L^{-1})^{\pm1}$ has a $\logq$-free representative; this applies to the edges, leaves and genus-$g$ graph terms built from the rows (\cref{def:mai-Fcf}).
\item The covariance $\Qr^{(i)}_{a,n}=\zeta^{i(a+1-n)}\Qr^{(0)}_{a,n}$ holds for the representatives, hence for the conifold expansions and for the hatted rows $\hat{\Qr}_a=s^{1-\nu_a}\Qr_a$, $z=s^4\zh$.
\end{enumerate}
\end{cor}

\begin{proof}
(i) The $\logq$-free subring of $\Nring$ (generated by $L^{\pm1}$, the $c_j$, the $i_j$, $c_0^{-1}$ and $\mathsf W^{-1}$) is stable under $\Dq$: $\Dq c_3$ is defined by $\PFfour c_0=0$, which does not involve $\logq$, and in the relation $\Conc(c_0,c_0\logq+i_0)=0$ that defines $\Dq i_2$ the coefficient of $\logq$ is $\sum_{a,b}c_a\Jet_{ab}c_b=0$, since $\Jet$ is antisymmetric (\cref{lem:gau-conc}~(iii)). The forward relations of \cref{prop:set-cyclicLR} for colour $0$, taken at $z^n$, give the stated recursion, whose coefficients are $\logq$-free elements of $\Nring$ (\cref{prop:gau-quintic}~(i)); as $\ev$ is a homomorphism of differential rings, the representatives evaluate to the row coefficients. For colour $i$ use \cref{lem:gau-colrows}. (ii) $\Delta_i^{1/2}=\zeta^{4i}C_0/L$ is represented by $\zeta^{4i}c_0L^{-1}$, a unit of $\Nring\otimes\bbQ(\zeta)$. (iii) The covariance holds for the representatives by construction; the conifold expansion is a $\bbQ(\zeta)$-linear ring homomorphism, and hatting multiplies the coefficient of $z^n$ in row $a$ by $s^{1-\nu_a+4n}$, independently of $i$.
\end{proof}

\begin{cor}[The colours of $\Dn$]
\label{cor:gau-colours}
For every $i\in\bbZ/5$, $\Dn_i(z):=\Dn_0(\zeta^{-i}z)$ satisfies, in $\bbQ(\zeta)\lser s\fps z$ and for the conifold expansions of the rows,
\[
\Qr^{(i)}\,\Dn_i(z)=\Pc\,\Qr^{\mathrm m,(i)}(v(u)),
\]
and it is the unique series in $\bbQ(\zeta)\fps z$ with this property. Hence, with $\Dn:=\diag(\Dn_0,\dots,\Dn_4)$ and the matrices of rows (columns indexed by the colours),
\begin{equation}
\label{eq:gau-R0}
\Qr\,\Dn=\Pc\,\Qr^{\mathrm m}(v(u)).
\end{equation}
Every $\Dn_i$ lies in $1+z\bbQ(\zeta)\fps z$, with $\coef{z^n}\Dn_i=\zeta^{-in}\coef{z^n}\Dn_0$; so $\Dn_0=1$ implies $\Dn=1$. In particular the expanded rows of every colour form a WKB solution of the quintic system with exponential $\re^{u_i/z}$.
\end{cor}

\begin{proof}
The ring $\cR_\zeta:=\bbQ(\zeta)\lser s\fps z$ is a power-series ring over a field, hence a domain; it contains $\Pc$ (with $u=s^5$), the model columns ($\varsigma^{-1}=s^{-1}\tau^{1/4}$) and the expanded rows. On $\cR_\zeta^5$ the colour rotation $\operatorname{Rot}_i$ of \cref{sec:gau-U} is additive, and $\operatorname{Rot}_i(PY)=(\operatorname{Rot}_iP)(\operatorname{Rot}_iY)$, $\operatorname{Rot}_i(Y\mathsf d(z))=\operatorname{Rot}_i(Y)\,\mathsf d(\zeta^{-i}z)$ (\cref{prop:gau-graded}~(iv)). By \cref{lem:gau-colmod}~(iii), $\operatorname{Rot}_i\Pc=\Pc$ (\cref{def:gau-can}), \eqref{eq:gau-R0zero} and \cref{cor:gau-rowsNN}~(iii),
\begin{align*}
\Pc\Qr^{\mathrm m,(i)}&=\Pc\operatorname{Rot}_i\Qr^{\mathrm m,(0)}=\operatorname{Rot}_i\bigl(\Pc\Qr^{\mathrm m,(0)}\bigr)\\
&=\operatorname{Rot}_i\bigl(\Qr^{(0)}\Dn_0(z)\bigr)=\operatorname{Rot}_i(\Qr^{(0)})\,\Dn_0(\zeta^{-i}z)=\Qr^{(i)}\,\Dn_0(\zeta^{-i}z).
\end{align*}
If $\Qr^{(i)}\mathsf d=\Qr^{(i)}\mathsf d'$ for scalar series $\mathsf d,\mathsf d'$, row $0$ gives $\Qr^{(i)}_0(\mathsf d-\mathsf d')=0$; the coefficient of $z^0$ of $\Qr^{(i)}_0$ is the expansion $\zeta^is^{-1}/\varpi_0\neq0$ of $\zeta^iL/C_0$, so $\Qr^{(i)}_0$ is not a zero divisor and $\mathsf d=\mathsf d'$. $\Dn_0\in1+z\bbQ(\zeta)\fps z$ by \cref{thm:gau-R0}~(iii), and $\coef{z^n}\Dn_i=\zeta^{-in}\coef{z^n}\Dn_0$ by the definition of $\Dn_i$. Finally $\Qr^{(i)}=\Pc\Qr^{\mathrm m,(i)}\Dn_i^{-1}$, with $\Dn_i^{-1}$ constant, is a WKB solution by \cref{thm:gau-R0}~(i).
\end{proof}

Gradedness of $\Pc$ is essential here: for the regular gauge $\Pc\cn(z)$ the series defined in the same way are $\Dn_i\cn(z)$, and they satisfy the colour relation for all $i$ if and only if $\cn\in\bbQ\fps{z^5}$. Once $\Dn_0$ is known to have rational coefficients (\cref{thm:con-glue}; see the remark after \cref{thm:gau-R0}), $\Dn_i$ for $i\neq0$ lies in $\bbQ\fps z$ if and only if $\Dn_0\in\bbQ\fps{z^5}$, since $\zeta^{-in}\notin\bbQ$ for $5\nmid n$.

\begin{defn}[The window $\Fg{g}$]
\label{def:gau-Fg}
For $g\geq2$, $\Fg{g}$ is the statement $\Dn_0\equiv1\pmod{z^{3g-2}}$; by \cref{cor:gau-colours} it implies $\Dn\equiv1\pmod{z^{3g-2}}$. It is the hypothesis of the comparison theorem (\cref{thm:flh-main}): at genus $g$ the leg-free graph sums involve the $R$-matrix only through its coefficients of $z^k$ with $k\leq3g-3$. By \cref{thm:gam-D}, $\Dn=1$, so $\Fg{g}$ holds for every $g\geq2$.
\end{defn}



